% ============================================================
%  R. Nielsen, "Den propædeutiske Logik" (1845)
%  "The Propaedeutic Logic"
%  Kjøbenhavn: Forlagt af Boghandler P. G. Philipsen;
%  Trykt i BiancoLunos Bogtrykkeri, 1845.
%  TRANSLATION (English) — from transcription.tex (the source of truth),
%  itself a diplomatic transcription of the Royal Danish Library scan.
%
%  Structure mirrors transcription.tex 1:1: Indledning (I., II.) — Første Deel
%  (Cap. 1 §§1–5, Cap. 2 §§6–8, Cap. 3 §§9–11) — Anden Deel (Indledning § 12,
%  Cap. 1 §§13–15, Cap. 2 §§16–18, Cap. 3 §§19–21). The §§ run continuously
%  1–21 across both Parts. The Første Deel Indledning's two sections are
%  numbered I./II., NOT §§; the Anden Deel's Indledning IS § 12 — that
%  asymmetry is in the book and is reproduced.
%
%  TERMINOLOGY (keep consistent; Nielsen's Hegelian-Danish school vocabulary):
%    Begreb = concept                 Forestilling = representation
%    Almeenforestilling = general representation
%    Anskuelse = intuition            Sandsning = sensation
%    Tænkning = thinking              Forstand(ig) = (the) understanding
%    Fornuft = reason                 Viden = knowledge
%    Erkjendelse = cognition          Indsigt = insight
%    Væren = being                    Væsen = essence
%    Virkelighed = actuality          Gyldighed = validity
%    Dom = judgment                   Slutning = inference
%    Beviis = proof                   Heelhed = wholeness/whole
%    Eenhed = unity                   Mangfoldighed = multiplicity
%    Ophævelse = annulment (Hegel's Aufhebung; "sublation" where the
%      dialectical double sense is clearly in play)
%    Aand = spirit                    Tilværelse = existence
%  The old id-est mark ɔ: → "i.e." (per the repo playbook). Danish „…“ →
%  English ``…''. \emph{} (Sperrsatz) carried over 1:1; antiqua \textit{}
%  for foreign-language matter carried over 1:1. Greek copied verbatim.
%  NO FOOTNOTES exist in this book; the \anm remarks are inline paragraphs.
%
%  Page-offset (per RESUME-NOTES): PDF = printed + 9.
% ============================================================
\documentclass[12pt]{book}
\usepackage[utf8]{inputenc}
\usepackage[T1]{fontenc}
\usepackage[english]{babel}
\usepackage[protrusion=true,expansion=true]{microtype}
\usepackage[margin=1.4in]{geometry}
\usepackage{setspace}
\usepackage{amsmath}
\usepackage{enumitem}
\usepackage{libertinus}
\usepackage{libertinust1math}
\usepackage{textalpha}   % polytonic Greek copied verbatim from the Danish
\usepackage{fancyhdr}
\usepackage[colorlinks=true,linkcolor=black,urlcolor=black,
  unicode=true,plainpages=false,pdfpagelabels]{hyperref}
\setlength{\headheight}{15pt}
\pagestyle{fancy}
\fancyhf{}
\fancyhead[LE,RO]{\thepage}
\fancyhead[RE]{\textit{The Propaedeutic Logic}}
\fancyhead[LO]{\textit{\leftmark}}
\setlength{\parindent}{1.5em}
\setlength{\parskip}{0pt}
\onehalfspacing

% ------------------------------------------------------------
%  STRUCTURAL MACROS — mirrored from transcription.tex so the two PDFs
%  render the book's four recurring display forms identically.
% ------------------------------------------------------------

% A Part opening: full-measure double rule, then the three-line head.
\newcommand{\deel}[2]{%
  \par\vspace{1.2\baselineskip}%
  \noindent\rule{\linewidth}{1.2pt}\par\vspace{2pt}%
  \noindent\rule{\linewidth}{0.4pt}\par\vspace{0.9\baselineskip}%
  \begin{center}
    {\large\bfseries The Propaedeutic Logic's\par}\vspace{0.35\baselineskip}
    {\bfseries #1\par}\vspace{0.35\baselineskip}
    {\Large\bfseries\emph{#2}\par}
  \end{center}
  \vspace{0.45\baselineskip}
  \begin{center}\rule{0.20\linewidth}{0.4pt}\end{center}
  \vspace{0.45\baselineskip}
  \phantomsection
  \addcontentsline{toc}{section}{\texorpdfstring{#1\quad #2}{#1 #2}}%
  \markboth{#1\ #2}{#1\ #2}%
}

% A Capitel head. As in the transcription: \capitel = plain argument
% (pp. 6, 35, 68), \capitelsp = letterspaced argument (pp. 114, 178, 247).
\newcommand{\capitelhead}[3]{%
  \par\vspace{0.9\baselineskip}%
  \begin{center}
    {\bfseries\emph{#1}\par}\vspace{0.25\baselineskip}
    {\large #3\par}
  \end{center}
  \vspace{0.3\baselineskip}
  \phantomsection
  \addcontentsline{toc}{subsection}{\texorpdfstring{#1\quad #2}{#1 #2}}%
}
\newcommand{\capitel}[2]{\capitelhead{#1}{#2}{#2}}
\newcommand{\capitelsp}[2]{\capitelhead{#1}{#2}{\emph{#2}}}

% A § head: centred "§ N.", then the centred bold argument.
\newcommand{\parag}[2]{%
  \par\vspace{0.9\baselineskip}%
  \begin{center}
    {\bfseries \S\,#1.\par}\vspace{0.2\baselineskip}
    {\bfseries #2\par}
  \end{center}
  \vspace{0.25\baselineskip}
  \phantomsection
  \addcontentsline{toc}{subsubsection}{\texorpdfstring{\S\,#1.\quad #2}{§ #1. #2}}%
}

% The Indledning's two sections, headed with a bare roman numeral (I./II.).
\newcommand{\romsec}[2]{%
  \par\vspace{0.9\baselineskip}%
  \begin{center}
    {\large\bfseries #1\par}\vspace{0.3\baselineskip}
    {\large\bfseries #2\par}
  \end{center}
  \vspace{0.3\baselineskip}
  \phantomsection
  \addcontentsline{toc}{subsection}{\texorpdfstring{#1\quad #2}{#1 #2}}%
}

% A top-level division head ("Introduction." — twice: p. 1 and p. 97).
\newcommand{\division}[1]{%
  \par\vspace{1.0\baselineskip}%
  \phantomsection
  \addcontentsline{toc}{section}{#1}%
  \markboth{#1}{#1}%
  \begin{center}{\Large\bfseries #1}\end{center}
  \begin{center}\rule{0.35\linewidth}{0.4pt}\end{center}
  \vspace{0.3\baselineskip}%
}

% An "Anm." remark: bold letterspaced lead-in, then an ordinary paragraph.
\newcommand{\anm}[1]{\emph{\textbf{Remark~#1.}}}

% The two tail ornaments, mirrored from the transcription:
% \divrule closes a top-level division; \secrule closes a §.
\newcommand{\divrule}{\begin{center}%
  \rule{0.27\linewidth}{0.9pt}\\[1.6pt]\rule{0.22\linewidth}{0.4pt}\end{center}}
\newcommand{\secrule}{\begin{center}%
  \rule{0.20\linewidth}{0.4pt}\\[1.3pt]\rule{0.20\linewidth}{0.4pt}\end{center}}

\begin{document}

\hypersetup{pdftitle={The Propaedeutic Logic},pdfauthor={Rasmus Nielsen}}

\frontmatter

% ============================================================
% TITLE PAGE — mirrors the 1845 title page (PDF 6).
% "Lie. theol." is the original's wrong-sort error for "Lic. theol.";
% normalised in translation.
% ============================================================
\begin{titlepage}
\centering
\vspace*{3.0cm}
{\Huge\bfseries The Propaedeutic Logic\par}
\vspace{1.6cm}
{\small by\par}
\vspace{0.7cm}
{\large\textit{Lic.\ theol.}\ {\LARGE\bfseries R. Nielsen,}\par}
\vspace{0.4cm}
{\large Professor of Philosophy.\par}
\vspace{1.8cm}
{\Large $\ast$\par}
\vspace{1.8cm}
{\Large\bfseries Copenhagen.\par}
\vspace{0.6cm}
{\emph{Published by the Bookseller P. G. Philipsen.}\par}
\vspace{0.35cm}
{\small\emph{Printed at BiancoLuno's Printing House.}\par}
\vspace{1.2cm}
{\large\textit{1845.}\par}
\vspace{0.8cm}
{\small Translated from the Danish.\par}
\end{titlepage}

% ============================================================
% CONTENTS — mirrors the original Indhold (PDF 8–9), which carries no page
% numbers against any entry. The generated \tableofcontents below carries
% this edition's own page numbers.
% ============================================================
\chapter*{Contents}
\markboth{Contents}{Contents}
\begin{center}\rule{0.35\linewidth}{0.4pt}\end{center}
{\small\setlength{\parindent}{0pt}\setlength{\parskip}{2pt}

\begin{center}\emph{Introduction.}\end{center}
\begin{center}I.\\ Formal and Speculative Logic.\end{center}
\begin{center}II.\\ The Aim of the Propaedeutic Logic.\end{center}

\begin{center}{\bfseries First Part.}\\ {\bfseries\emph{The Doctrine of the Subjective Concept.}}\end{center}

\begin{center}{\bfseries\emph{First Chapter.}}\\ The Formation of the Subjective Concept.\end{center}
\S\,1.\quad The General Representation.\par
\S\,2.\quad The Clear and Distinct Intuition.\par
\S\,3.\quad Insight, or the Subjective Concept.\par
\S\,4.\quad The Apriority of the Concept.\par
\S\,5.\quad The Totality of the Concept.\par

\begin{center}{\bfseries\emph{Second Chapter.}}\\ The Presentation of the Subjective Concept: the Logical Judgment.\end{center}
\S\,6.\quad The Logical Form of the Judgment.\par
\S\,7.\quad The Logical Content of the Judgment.\par
\S\,8.\quad The Logical Totality of the Judgment.\par

\begin{center}{\bfseries\emph{Third Chapter.}}\\ The Validity of the Subjective Concept: Inference and Proof.\end{center}
\S\,9.\quad The Logical Requisites of Inference.\par
\S\,10.\quad The Formal Validity of Inference.\par
\S\,11.\quad The Real Validity of Inference: Proof and Method.\par

\begin{center}{\bfseries Second Part.}\\ {\bfseries\emph{The Doctrine of the Objective Concept.}}\end{center}
\begin{center}\emph{Introduction.}\end{center}
\S\,12.\quad The Dialectical-Speculative Method.\par

\begin{center}{\bfseries\emph{First Chapter.}}\\ The Formation of the Objective Concept.\end{center}
\S\,13.\quad Being and Essence.\par
\S\,14.\quad Essence and Actuality.\par
\S\,15.\quad Actuality and the Concept.\par

\begin{center}{\bfseries\emph{Second Chapter.}}\\ The Freedom of the Objective Concept.\end{center}
\S\,16.\quad Ideality and Reality.\par
\S\,17.\quad Subjectivity and Objectivity.\par
\S\,18.\quad Teleology.\par

\begin{center}{\bfseries\emph{Third Chapter.}}\\ The Fullness of the Objective Concept: the Logic of the Idea.\end{center}
\S\,19.\quad The Originality of the Idea.\par
\S\,20.\quad The Interaction of the Ideas.\par
\S\,21.\quad Idea and God.\par

}
\begin{center}\rule{0.35\linewidth}{0.4pt}\end{center}

\cleardoublepage
\tableofcontents

\mainmatter

% ============================================================
%  BODY — Danish printed pp. 1–283 (PDF 10–292)
% ============================================================

% ---- Introduction, printed pp. 1--5 ----
\division{Introduction.}

\romsec{I.}{Formal and Speculative Logic.}

% Danish p. 1 opens with a raised initial "V"; only the absent indent is
% reproduced, as in the transcription.
\noindent By reflecting upon its own thinking, the subject is led to make
clear to itself the various operations of thought and their mutual
connection. The first requirement is that thinking be consistent and follow
its own laws without any regard either to the psychical peculiarity of the
thinking subject or to the character of the objects to which it might
otherwise be applied. Formal logic therefore treats thinking as a given
phenomenon of spirit, without asking after its origin or its objective
validity. This question is nonetheless so momentous that the whole direction
and destiny of philosophy depends upon the way in which it is answered. The
sensualists teach that the soul (\textit{tabula rasa}) receives its original
content from the sensory impressions of the surrounding world (impressions),
and end by making thinking a mere shadow of experience. Subjective idealism,
on the contrary, vindicates the soul's infinite self-activity, and conceives
the world of experience as the representing I's own unconscious product. Each
of these theories misjudges human thought. Sensualism undervalues thinking,
doubts the universal validity of the laws of thought, and can for that reason
find no truth in the actual world; idealism overvalues subjective thinking,
denies the world's true actuality, and ends in formal emptiness. Both
one-sidednesses are overcome in speculative logic, which shows that our
rational thinking mirrors the reason revealed in the universe.

\anm{1} By making sensation the source of cognition, sensualism was bound to
lead to skepticism; for sensation refers only to the particular and
contingent, not to the universal and necessary; experience teaches merely
\emph{that} something happens, but not \emph{why} it happens; it shows that
the one follows upon the other, but does not prove that the one is a
consequence of the other; on such presuppositions it was accordingly
consistent that Hume contested the concept of causality.

\anm{2} Against Hume's theory Kant raised his opposition. The fundamental
thought of the Kantian critique of reason is this, that our pure sensory
intuitions (time and space), categories of the understanding (unity and
multiplicity, cause and effect, etc.) and ideas of reason (the soul, the
world, God) cannot have been abstracted from a world of experience lying
outside consciousness, since they themselves constitute the necessary
presuppositions without which no experience is possible. But while Kant
defended the subjective universal validity of the a priori forms, he denied
them, on the other hand, all objective significance. The whole world of
experience is consequently to be regarded as a product of two factors, namely
of the a priori forms and of an unknown object (``das Ding an sich''), whose
pure essence no one can cognize, because everyone who attempts to cognize it
thereby at once takes it up into the subjective medium of his own
consciousness.

\anm{3} The Kantian ``Ding an sich'' was contested by the elder Fichte. It
is, according to Fichte's view, an absurdity to assume a ``Ding an sich,''
since such a thing can neither be sensed nor thought. To an object that was
supposed to lie wholly outside consciousness the subject cannot ascribe
reality without falling into contradiction with itself; for if anyone says
that he knows of the existence of such an object, then he has already drawn
this existence into the circle of consciousness; he then says: I know
something of an existence of which I nevertheless know nothing at all.
``Das Ding an sich'' was the last remnant of objectivity that the critical
idealism left standing. By removing this unknown quantity, idealism became
wholly subjective, and now concentrated itself in the proposition that the
whole world is brought forth by the I. In order to understand this bold
assertion, we must determine more closely the I's essence, its productivity,
and its freedom.

\begin{itemize}
\item[a)] The I of which idealism speaks is not the empirical I, which finds
  itself immediately determined by the various expressions and states of the
  life of consciousness (I speak, grieve, rejoice, etc.); but the pure
  I~=~I, the I that in everything thinks itself, without regard to
  particular limitations and alterations.
\item[b)] The productivity whereby the I brings forth the whole actual
  surrounding world is no loose play of fancy, nor an arbitrary imagination
  which the I could alter or arrest at its pleasure, but an act of
  necessity, whereby it presupposes those barriers whose infinite annulment
  conditions its freedom.
\item[c)] The I's production of the world goes on unconsciously in its own
  obscure ground of life, and belongs, so to speak, to its nocturnal sphere;
  hence the universe comes forward with so imposing a power that the subject
  takes its own vision of the world for an outward, objective actuality.
  This popular view, however, is only a dream, which vanishes when
  philosophical thinking awakens. The transition from the standpoint of
  representation to the standpoint of philosophy is thus a transition from
  illusion to truth, from unfreedom to freedom.
\end{itemize}

\anm{4} Speculative logic has in recent times been founded by Hegel.

\romsec{II.}{The Aim of the Propaedeutic Logic.}

It is thus a matter of fact that besides the so-called thinking of the
understanding, which finds its immediate application in practical life and
its higher cultivation in mathematics and the empirical sciences, there is
yet another kind of thinking, namely the dialectical, which conditions all
higher knowledge, and therefore has its proper home in speculative
philosophy. Although both these main directions of thinking have their
origin in one and the same source, the eternal reason, they nevertheless
form, for the superficial view, so striking an opposition that it often
seems as though truth were at strife with itself. This opposition shows
itself not merely in a divergent terminology, but even in the way in which
the laws of thought themselves are determined and presented, so that one may
be very well practised in the thinking of the understanding without yet
finding oneself able either to follow the course of philosophy or to
appropriate its results. If, then, the propaedeutic logic is to answer to
its true destination, if it is to orient us in the whole realm of thinking,
it must necessarily engage with both of the opposed directions of thought.
By keeping to formal thinking alone it would evidently give too little; by
losing itself in the solution of speculative problems it would overstep its
propaedeutic boundary: it therefore occupies a mediating place between
formal and speculative logic.

\anm{1} In order to accomplish its task, the propaedeutic logic must give a
comparative presentation of the formal and the speculative operations of
thought. If this presentation is to succeed, it must be shown step by step
how a determinate thought-content, after first being regarded with the eye
of the understanding, changes its appearance when viewed through the glass
of dialectic. \emph{The Doctrine of the Subjective Concept}, which
constitutes the first part of the logic, therefore lets thinking seek its
points of departure, step by step, in the firmer soil of the sense-world and
of representation, in order thence to rise into the freer air-current of
speculative cognition. By this movement from below upward we are everywhere
led from experience into philosophy, from the circle of thought we call our
own into that world of reason which reveals the divine Thinker. The exercise
of thought and the insight into the essence of the concept that are hereby
won introduce, in the next place, a more methodical development of the
universal forms and principles of thinking and of existence, which is
presented in the second part of the logic, i.e.\ \emph{in the Doctrine of
the Objective Concept}. But just as the first main division constantly
points forward to a higher development of the concept, so must the latter in
turn everywhere point back to the concepts' empirical existence, and cast
light upon the actual world.

\divrule

% ---- First Part, printed pp. 6--96 ----
% The body head prints "Første Deel:" with a colon (the Indhold has a period);
% the colon is what p. 6 sets and is mirrored here.
\deel{First Part:}{The Doctrine of the Subjective Concept.}
\capitel{First Chapter.}{The Formation of the Subjective Concept.}
\parag{1}{The General Representation.}

% Danish p. 6 opens flush left with a raised initial "L"; only the absent
% indent is reproduced.
\noindent Likeness and unlikeness are the simplest forms in which the objects
of the surrounding world can be presented to thinking consideration. The
operation of thought that consists in bringing out and holding fast the like
in unlike objects is an \emph{abstraction}, and the result to which it leads
a \emph{general representation}, which may be characterized as a sense-image,
a thought-image, and the fore-image of the concept.

Although two existing things cannot, in one and the same respect, be both
like and unlike each other, it is nevertheless impossible to think likeness
except by the help of unlikeness; the more sharply these categories are
separated, the more clearly it appears that they pass over into each other:
their boundary is consequently dialectical.

\anm{1} Several objects that have an essential likeness to one another we
reduce in language to a single common designation, and thereby show that we
overlook their accidental unlikeness. With language, abstraction begins at
once; as the child learns to speak, it learns at the same time to abstract.
The likeness of objects is observed by a comparison that is at first just as
immediate as sensation itself, and abstraction proceeds, so far, quite
unconsciously; but since likeness itself is no sensory object that could be
removed, by outward separation, from the unlikenesses in which it finds its
expression, abstraction shows itself to be an ideal act, whereby the subject
itself takes the universal for the essential and lets the accidental
modifications go --- hence a beginning act of thinking, which needs only to
be brought to consciousness in order that thinking proper may begin.

\anm{2} To see an object is, in other words, to form for oneself a
representation corresponding to determinate sensory impressions. That the
subject itself forms the representation is evident from the fact that the
image of the object can be reproduced after the object itself has
disappeared. By sensing a series of different objects there is developed in
consciousness a circle of different images and representations; by the
application of comparison and abstraction a general representation is
obtained. The general representation is in the beginning quite immediate.
One represents the universal to oneself by setting up a single object as
exemplar. Universality thus coincides with the image in which it is
intuited, and is consequently to be regarded as a \emph{sense-image}. But a
closer examination shows that every exemplar that is to represent pure
universality gives us both too much and too little. Besides the universal,
the exemplar expresses also the peculiarity of a determinate object, and so
contains too much; on the other side it shows universality only in a
particular shape, and for that reason gives us too little. The subject, which
now attempts to form a determinate representation of the universal, obtains
then only an unsurveyable series of exemplars, each of which, taken by
itself, is equally adequate and equally inadequate. The consequence is that
the general representation, encumbered with sensuousness, vanishes like a
phantom-image, and only the thought of the universal, which is itself a
thought, remains behind. By passing over into thinking, the general image
becomes, properly speaking, imageless. Thinking is negative insofar as the
universal is at the same time the invisible, which cannot be represented or
held fast in any image. Abstracting thought sets the universal free from all
sensory admixtures; but when it is further considered that the universal
never exists for itself, but is abstracted from the existing objects, then
thinking must in turn annul the abstraction, and the eye of the
understanding must behold the universal in its single exemplars and
examples. Free thought hovers unceasingly between the imageless and the
imaged, whereby the general image is transformed into a
\emph{thought-image}. Accordingly it encloses \emph{and comprehends} all the
single objects, and in this sense is to be posited as \emph{the fore-image
of the concept}.

\anm{3} When the universal is held fast by abstraction, it is called an
abstractum; when on the contrary it is intuited in a determinate exemplar,
we have it \textit{in concreto}. Yet the relation between the abstract and
the concrete is determined otherwise in formal than in speculative logic.

\anm{4} As regards the relation between pure likeness and unlikeness, it is
in the first place evident that they mutually exclude each other, or, in
other words, that they bound each other. But this does not settle the
matter; for every negative bounding of thought must also be converted into a
positive determination of thought. What is the like? Answer: the like is
unlike the unlike, and thereby like itself. In this lies that the thought of
likeness not only falls outside the thought of unlikeness, but also
coincides with it. He who thinks directly of the like thereby thinks
indirectly of the unlike. Unlikeness is thus the category whereby likeness
is determined. The boundary between pure likeness and unlikeness can never
stand still; hence it is said, in the language of dialectic, that they pass
over into each other. That this dialectic is not arbitrary is proved by the
history of philosophy itself. For by making the boundary between the like
and the unlike immovable and undialectical, the Eleatics could prove that
all becoming was unthinkable. If becoming is to be thought, they said, then
it must be either the like that comes from the unlike, or the like that
arises from the like. In the latter case no becoming can take place, for
since the like is = the like, there is no transition here; in the former
case something would have to come to be out of nothing, for in the unlike
the like cannot be contained at all. The Eleatic abstraction betrays its
one-sidedness in that it will not let the one concept integrate the other.

\parag{2}{The Clear and Distinct Intuition.}

When the object of thinking consideration shows itself as a \emph{unity} in
a multiplicity, the subjective activity of thought consists: a) in an
\emph{analysis} or ideal singling-out, whereby the manifold marks and
determinations that are to be distinguished with respect to the object are
set apart from one another; b) in an ideal gathering-together or
\emph{synthesis}, whereby the singled-out plurality is led back to its
unity. The result won by this double operation of thought is a \emph{clear
and distinct intuition}. By virtue of the intuition's clearness, the
objectual unity is singled out from the infinite multiplicity of existence
and posited as determinate singleness; by the intuition's being made
distinct, the single object is itself determined by a multiplicity of single
features, which are comprised within its own unity. Without an analysis and
synthesis, albeit an unconscious one, no determinate sensation or experience
is possible; but the analytic and synthetic activity of thinking is from the
beginning so wholly lost in sensation that the clear and distinct intuition
itself bears the immediate stamp of sensuousness. Nevertheless the concept
is here too present as fore-image; for analysis and synthesis cannot be
applied except to relations; but a relation is itself no sensory object.

When the immediate categories of relation --- unity and multiplicity --- are
themselves made the object of a thinking inquiry, a dialectical problem
arises.

\anm{1} Insofar as the universal is the common form in whose likeness the
mutual unlikeness of the single exemplars vanishes, it shows itself as the
unity under which the objects can be subsumed; insofar as it also finds its
concrete expression in each exemplar by itself, its likeness is itself
further determined by unlikeness, maintains itself in a repetition, and
passes over into plurality. Hence also most \textit{nomina apellativa} have
% "apellativa" with a single p is what the Danish page sets; carried as
% transcribed, per the transcription's note.
both a singular and a plural.

\anm{2} When that which is apprehended refers not to the sense of sight but
to that of hearing, there can indeed be no talk of intuition, but there can
always be talk of a clear and distinct apprehension.

\anm{3} When one speaks of an eye of the understanding, or of a
thought-vision, these expressions must not be understood merely figuratively;
for the transition from sensation to thinking cannot take place by a leap.
As a thinking lies at the ground of all sensation, so an ideal sensing also
lies at the ground of free thinking; as there is a sensory intuition, so
there is also an intellectual intuition, and just as one may speak of a
physical blindness, one may also speak of a mental blindness. But when it is
now to be stated with determinateness what that ideal is toward which
thought's gaze turns, there at once arises the difficulty that thought never
dwells upon anything that can be set out and pointed to as a determinate
example or a resting shape. While sense settles down to rest with a bounded
object, thought is on the contrary in eternal unrest, and oversteps every
boundary. As the sensing eye dwells on the one thing, it has at the same
time turned itself away from the other; when thought, on the contrary, is to
hold fast one determination, it always takes another to its aid. Hence the
clear thought-vision is properly a double vision: a \emph{regard}
(\textit{Hensyn}). What the \emph{regard} is from the subjective point of
view, the \emph{relation} is from the objective; therefore thinking always
aims at finding the determinate relations. In consequence, analysis and
synthesis acquire a quite different significance for intellectual than for
sensory intuition. For sense, the intuition is clear when the object in its
pure singleness comes forward with determinate outlines; but thinking does
not forget that the singled-out object is itself one among several, that its
singleness can therefore be determined only with regard to other single
things, and that it must accordingly be posited as a \emph{separate} object.
From the sensory side, the intuition is made distinct when the plurality
contained in the object's unity comes forward as a sum of single features;
for thought, on the contrary, the intuition can be made distinct only by the
given single features being determined with regard to one another, and
posited as members of separate relations.

\anm{4} The relation of unity to plurality is determined by merely
representational thinking thus: that something which in one respect can be
regarded as a plurality can in another respect be posited as a unity; but
dialectical thinking is not content with so external a joining of opposed
regards; it asks whether unity, when clearly and distinctly thought through,
does not itself determine itself as a plurality. In Plato's dialogue, the
Parmenides, Socrates has hit upon the problem, which he propounds roughly
thus: That I, Socrates, am in one respect a unity, in another respect a
multiplicity, is not so wonderful; but if anyone could prove to me that ἕν
was πολλὰ and πολλὰ ἕν, then I should marvel greatly (ἀγαίμην ἂν θαυμαστως,
ὦ Ζήνων). The Eleatics (Zeno) demonstrated the contradiction that lies in
the category of plurality. ``If it is assumed that several things exist,
then their number must be both determinate and indeterminate. The things
must be thought to make up a determinate number, for they must be just as
many as they are (neither more nor fewer); but on the other side their
number must also be thought to be indeterminate; for between the separated
things there are yet other things, and outside the number of determinate
things again a number of determinate things; consequently the number of
things becomes infinite, indeterminable, and thereby in itself
indeterminate. Since it is now a contradiction to think the number of things
both as determinate and as indeterminate, plurality is assumed to be
unthinkable.'' The truth in this is that abstract analysis can go on to
infinity: every plurality consists of unities, each of these unities is
again to be thought as a plurality with its unities, and so forth. Abstract
synthesis likewise goes on to infinity; for the unity in which a plurality
is gathered is again a member in a series of several similar unities; if
these are now bound together by a new synthesis, we have once more a unity
in a new series, and so on. In this respect we must therefore grant the
Eleatics their point when they maintain that the number of things will
become infinite if a plurality is assumed at all. But the thought of such an
infinity stands in no insoluble contradiction with the proposition that the
number must also be determinate. For when synthesis is everywhere joined
with analysis, every unity becomes determined by \emph{its} plurality, and
the number of things in consequence \emph{infinitely determinate}.

\anm{5} If the relation between unity and plurality is considered purely
quantitatively, unity is posited as continuity, plurality as discreteness.
Every line represents such a continuity, whose discreteness is again to be
sought in the line's innumerable points.

\parag{3}{Insight, or the Subjective Concept.}

If an existing object is also to be an object for a comprehending thinking,
its content must let itself be apprehended in the form of \emph{wholeness
and parts}. As the whole and its parts are seen mutually to condition and
determine each other, questions are thereby called forth about
presupposition and consequence, about ground and consequent, and thinking,
which now renders to itself an account of the whole complex of determinate
relations, shows itself as a \emph{reflection}, whereby the subject comes to
a \emph{concept of the thing}, or to an \emph{insight} into the essence and
character of things.

The contradiction that lies in the concept of totality itself was brought
forward early by the Greek philosophers, just as it also plays an important
role throughout the whole history of philosophy, especially with respect to
the question of the infinite divisibility of matter. As a contradiction
between the simple and the composite it was reduced by Kant to an antinomy,
which in turn found its dialectical resolution in the Hegelian logic.

\anm{1} When unity is posited as wholeness, it is thereby expressed that it
embraces a multiplicity of single elements, which make up its parts. For
sensory consideration the whole shows itself merely as the aggregate of the
parts; for intellectual apprehension, on the contrary, the whole is to be
determined as the \emph{general} connection of the parts, and the parts as
members in the whole's \emph{separate} connection. When the whole is given
as presupposition, the character and mutual relation of the parts follow
with necessary consequence. Hence it is said in Professor \emph{Sibbern's}
explicative logic that by cognizing the whole's \emph{mode of existence} one
comes to see into the parts' \emph{ground of existence} (\textit{ratio
essendi}). For when something is posited directly in such wise that
something else is thereby posited indirectly, the former is called the
ground, the latter the consequent; therefore the ground is always the
primitive member of the relation, and the consequent its consecutive member.

\anm{2} From the determinate relation between wholeness and parts it is
evident that the true thought-vision is a double vision. As the whole is
seen immediately, the parts come mediately into view, and conversely. The
parts and the whole mirror or reflect one another; hence free,
self-conscious thinking has also received the name of \emph{reflection}.

\begin{itemize}
\item[a)] The word reflection is no doubt borrowed proximately from optics;
  but the natural light is also a copy of the light of thought. There is the
  same relation between thinking's opposed members as between ``two mirrors,
  of which the one reflects the other's image, and in this reflection
  reflects its own'' (cfr. Professor Heiberg: Grundtræk til Philosophiens
  Philosophie, Pg.\ 34).
\item[b)] Reflection is also the expression for the \emph{deliberating and
  testing} thinking. This signification of the word is derived from the
  foregoing. Only when the one thought is seen in the light of the other
  does it become evident how far they can subsist with each other or not.
\item[c)] Finally, reflection can also designate the \emph{reasoning and
  grounding} thinking; for when two thoughts strive against each other, it
  becomes a question which of them the subject is to give up, and which to
  hold fast; this can be decided only by a determinate regard, or by a
  presupposition that contains the determining ground. The thought of man's
  freedom, for example, cannot be reconciled with the thought of an
  inescapable fate: the thinking subject must now hold fast either freedom
  or fate; if regard is had to conscience and imputability, then the thought
  of an irresistible fate must be given up; if, on the contrary, regard is
  had merely to a mechanical necessity of nature, the thought of freedom
  must vanish. Which regard the subject will take in a given case may in
  turn depend upon other regards, and so on to infinity. Since every regard
  contains a partial ground, the thinking that is lost in external regards
  leads to a reasoning of grounds. Insofar as the actual surrounding world
  shows us things in their many different combinations and entanglements,
  reasoning thought becomes necessary in all empirical inquiries. But
  reasoning thought is nevertheless of finite nature; if it attempts to tear
  itself loose from given experience, it must end by adducing grounds
  \textit{pro} and \textit{contra}, without ever reaching the highest and
  last ground. Only when all outer regards vanish in one fundamental view
  does reflecting thought take on a philosophical character, and find the
  infinite presupposition from which all finite relations can be deduced.
  From the standpoint of finitude, thinking is thorough when it apprehends
  the matter in its many different relations, the remoter as well as the
  nearer; from the standpoint of infinity it is thorough only when it seeks
  and finds the absolute ground. As it is a philosophical affectation to
  want to apply the infinite forms of dialectic to the finite content of
  everyday life, or to disdain a reasoning of grounds where what is in
  question is an understanding critique or an orientation in practical
  affairs, so on the other hand it betrays a scientific narrowness to want
  to make one's way through with a merely reasoning thinking when what is at
  stake is a cognition of the eternal truths.
\end{itemize}

\anm{3} The result of reflection carried through is \emph{the concept}, as
such, or the proper thought-result. The concept is thus a thought-whole, or
a fundamental thought, that lets itself appear in a multiplicity of separate
thought-determinations and their mutual connection. If it is now asked where
the concept exists, it may in the first place be said to have its existence
in the thinking subject's consciousness. The subject sees the object in its
different relations, abstracts from it the determinate relations, sunders
(analysis) and unites (synthesis) the partial thoughts in their fundamental
thought (reflection), and is therefore said to form for itself a
\emph{concept of the thing}. But insofar as it is the existing object from
which the subjective concept is abstracted, one can also speak of the
object's own concept, or of the thought-whole expressed in the object. The
difference between the object and its concept is, properly speaking, a
difference between its intellectual and its sensory side, which is more
closely to be determined as an opposition between the thing's \emph{inner}
and its \emph{outer}. The outer is for sense, the inner for thought; hence
the subjective \emph{concept of the thing} is also said to be an
\emph{insight into the thing} (ἐπιστήμη --- νοητόν).

\anm{4} In order to form a concept of those phenomena that come into view
piecemeal, without the whole being immediately surveyable, it often becomes
necessary to imagine for oneself an intellectual whole, or to set up a
\emph{hypothesis}. The cognition to which the hypothesis leads is, however,
only finite and relative. Thought hovers between experience and reflection.
In empirical respect the hypothesis is insufficient; for the parts (the
outer) \emph{condition} but do not \emph{ground} the whole (the inner), and
in intellectual respect the insight rests upon a reasoning of grounds that
lacks the sufficient ground. The hypothesis therefore yields only a
probability, which must be raised to certainty either by a new experience,
which then introduces a new inquiry of thought, or by a deeper insight into
the connection's essential ground (the Copernican hypothesis, Newton). What
the hypothesis is in relation to the whole, \emph{the conjecture} is in
relation to a single member missing in the whole. (Jfr.\ Sibberns Logik
\S\,23.)

\anm{5} Representational thinking determines the relation between the whole
and the parts thus: the whole is posited = the sum of the parts, but greater
than any of its single parts. This way of seeing has its full validity so
long as only finite magnitudes are in question; but it is wholly
insufficient as soon as the inquiries are carried over into the territory of
infinity. The spiritualists (Berkeley), who denied the being of matter,
appealed to the contradiction that lies in the thought of an infinite
divisibility. Matter, according to their contention, cannot exist at all;
for since it is thought to consist of infinitely many infinitely small
parts, every part must become = the whole. Suppose, for example,
\textit{a} (the part) $<$ \textit{A} (the whole); by multiplying by $\infty$
on both sides one obtains $\infty$ \textit{a} = $\infty$ \textit{A}, and by
again dividing both factors by the same magnitude, namely $\infty$, there
results: \textit{a} = \textit{A}. The mathematicians' objection, that such
an operation is absurd, cannot refute the spiritualists' theory, for it
rests precisely upon the admitted absurdity. Equally insufficient is the
common remark that in the infinite itself no boundary can be set, and that
greatness and smallness must therefore flow together in the fog of infinity;
for the task is precisely to let the light of thought disperse this fog, or,
in other words, to cognize how the finite and the infinite mutually
determine each other. --- Still more determinately did the problem come
forward in the Kantian critique of reason, where it was reduced to an
antinomy lying in the concept of substance:

\emph{Thesis: Every composite substance must consist of simple parts.} If
not, then every part would in turn be composite, and the resolution of the
composition would have to be continued until at last the substance came to
consist of nothing.

\emph{Antithesis: There exists no simple substance.} For no composition is
possible except in space; now if there were a substance consisting of
simple parts, then space itself would have to consist of indivisible parts;
but space can be divided to infinity.

The truth in the Kantian antinomy is, as Hegel has also shown, this: that
the thought of a whole always carries with it the thought of parts. When a
certain determinate part is separated from its whole, it is itself posited
as a whole. The same thing that from one point of view must be regarded as a
whole must from another point of view be determined as a part; the relation
turns itself about unceasingly, and every finite boundary is annulled. But
each time a determinate whole is posited, it enters into a determinate
opposition to its parts, so that the finite side of the boundary also comes
into view. Everything then depends solely upon the relation itself being
more closely determined. The relation between wholeness and parts has a
certain likeness to the relation between unity and plurality, continuity and
discreteness; nevertheless an essential difference is here to be observed.
For when continuity is posited, discreteness has vanished; and when
discreteness is posited, continuity is broken off. Continuity and
discreteness can be exchanged with each other to infinity, but they do not,
properly speaking, reflect each other. This, on the contrary, is the case
with the whole and the parts; when the whole is posited, the parts are
posited too, and conversely. Hence this relation is called a \emph{relation
of reflection} or an \emph{essential relation}.

\anm{6} At its lowest stage, wholeness is a mechanical compositum, which
shows us the relation from its finite side. On a higher stage stands the
organic whole; hence there is also said to be an infinite relation between
life and organism. In its highest development, wholeness is raised to a free
totality, a reflected \emph{singleness}, in which the concept coincides
entirely with its object. Thus the whole I is present in each of its
self-conscious expressions of life, and yet it is a singleness that cannot
be divided. As sensation \emph{begins} immediately with the single, so
thinking \emph{ends} with the reflected singleness.

\parag{4}{The Apriority of the Concept.}

Although the subjective concept in its original development is conditioned
by experience, it nevertheless reflects its own essential determinations
with such independence and consistency (concept-analysis) that it can by no
means be a product or empty reflection of experience. Just as little can the
world of experience be regarded as a product of the subjective I, although
all experience is conditioned by a priori thinking. Hence formal logic also
makes a distinction between \emph{a priori and a posteriori concepts}. The a
priori concepts do not belong merely to the thinking subject, as though they
formed an impassable gulf between subjective cognition and objective truth,
but express those \emph{general} fundamental forms which do not presuppose
special experience, inasmuch as they everywhere find their application to
the existing objects and their different relations. On the other side, the a
posteriori concepts, which are abstracted from the objects met with in
experience, must not be thought to stand indifferent to the subjective laws
of thought, for thereby they would cease to be true concepts; they are much
rather to be regarded as the \emph{separate} forms in which a priori
thinking recognizes its own peculiar content (normal-concepts).

The infinite interaction between subjective thinking and objective reason
presupposes a creative freedom, as the absolute prius, from which all
thinking and being have their common origin.

\anm{1} Professor Sibbern, in his explicative logic (2nd edition Pag.\
190--200), has distinguished between the \emph{pursuing}, the
\emph{general}, and the \emph{pure concept-analysis}. The \emph{pursuing}
analysis carries the concept-determinations found in the object out beyond
the empirical boundary, by making distinct the consequences lying within
them. The \emph{general} concept-analysis finds the concept in the living
word, separates the word's original meaning from the derived, shows its
kinship with other words (synonymics), thereby discovers the essential
requisites of the concept sought, and is thus able not merely to determine
but even to correct linguistic usage. The \emph{pure} concept-analysis makes
distinct thinking's own categories, and plays a chief role in the objective
logic. In all these kinds of concept-analysis it appears that thinking goes
out beyond the given, so that the sensualists are wrong in denying its
independence.

\anm{2} The critical idealism, which indeed concedes to thinking a
subjective apriority, but denies its power to cognize objective truth, and
in proof of this appeals to the antinomies, finds its refutation in the
dialectical resolution of these antinomies. On the other hand the subjective
idealism, which refers the whole objective world-content to the I's infinite
apriority, cannot be refuted by the fact that the development of thinking is
conditioned by experience, since it is not the world of experience's
necessity but its independence that it denies.

\anm{3} The Cartesian principle: \textit{Cogito, ergo sum}, has found its
consistent development in the Fichtean ``Wissenschaftslehre.'' The absolute
subject, I, can doubt everything, only not its own thinking and being. The
proposition: I \emph{think}, falls immediately together with the
proposition: I \emph{am}. The subject cannot think without \emph{being}
thinking, and though it should doubt the validity of all single thoughts, it
still cannot by any doubt whatever annul the certainty of thinking in
general, since in order to doubt all thinking it must necessarily think, and
consequently rely upon thinking in the same instant in which it doubts
thinking's validity. Whatever does not have this incontestable certainty
must fall outside science proper. Since science cannot be an aggregate of
loose representations, since one fundamental truth must run through all
single truths, all scientific propositions must also let themselves be led
back to one irrefragable first principle, and from this in turn be deduced.
If now the I has found this general first principle in its own thinking and
being, then it has at the same time found the source of truth in itself, and
philosophy, the pure science, would then according to this theory be in a
position to comprehend the whole of existence in virtue of the supreme first
principle, or, in other words, to construct its world \textit{a priori}. The
truth in this is that scientific thinking can acknowledge nothing but what
the consistency of thought carries with it, and that the necessity with
which truth compels from the thinking subject its acknowledgment is the
purest touchstone of the certain and the uncertain. But the subjective
idealism's constructing \textit{a priori} carries with it nonetheless a
one-sidedness that shows itself in several ways (Jacobi's
knitting-stocking). The subject, which always has its gaze turned toward the
formal I-hood, loses sight of the actual world and its concrete content.
Either philosophical knowledge must then be restricted to a small number of
abstract truths, or else the philosopher who undertakes to construct nature
and history \textit{a priori} will easily be convicted of having
insinuated empirical presuppositions contrary to the principle set up.
Unsuccessful attempts to construct \textit{a priori} what can be cognized
only in experience's a posteriori medium have contributed above all else to
bringing philosophy into discredit.

\anm{4} The distinction between a priori and a posteriori concepts refers,
according to formal logic, to their source or origin. The simplest form in
which such a distinction can be held fast is this, that the a priori
concepts are said to be \emph{prior to} the things to which they are
\emph{applied}; the a posteriori, on the contrary, to follow \emph{after the
things} from which they are \emph{abstracted}, so that the former have their
origin from thinking itself, the latter from experience. But more closely
considered, this opposition cannot maintain itself. It has already appeared
that the concept of a thing is not to be sought in its accidental outer, but
in its inner, i.e.\ in its intelligible essentiality. That thinking
abstracts its empirical concepts from the objects means then, in other
words, that in their determinate essence it recognizes its own original
essentiality. Experience is thus no proper source of concepts, but only the
means whereby the subject's general thinking is enriched with a
thought-content concrete and distinctly stamped in itself. The concept,
which with respect to our apprehension may be a posteriori, insofar as a
special experience has gone before, is yet with respect to the object itself
a priori, since a thing's inner and outer always stand related as the
primitive (\textit{dignitate prius}) to the consecutive (\textit{dignitate
posterius}). The a priori interaction between the thing's concept and our
comprehending thinking comes into view in the development of the so-called
\emph{normal-concepts} (jfr.\ Sibbern \S\,47). The normal-concept begins
\textit{a posteriori}, insofar as its general ground-lines are first
discovered in an empirically given object; but it shows itself at the same
time from an a priori side, insofar as it at once becomes evident that the
actual object does not answer to its demands. This disproportion then leads
the subject, thinking through the concept's consequences, to imagine or
construct for itself an object so perfect that the intimated concept finds
therein its adequate expression. The fruitful application that can in
actuality be made of such normative concepts (for example in
crystallography) proves clearly that thorough thinking has found the
original meaning of things, or, what is the same, that the a priori demands
that guide the thinker also govern actuality. All so-called a posteriori
concepts, when they are not confused with mere representations or sensory
intuitions, are properly to be regarded as special appearances of the
universal reason, whose apriority is as objective as it is subjective. With
respect to the method of our comprehension, on the other hand, the
distinction between a priori and a posteriori concepts may so far have its
significance, as it reminds the thinking subject of its finitude and
intellectual limitation. Insofar as the concrete concepts are won \textit{a
posteriori}, experience must go hand in hand with comprehending thinking;
insofar as the a posteriori concepts are the a priori thinking's own types,
experience's sporadic content must be purified in the pure, universal
dialectic. From the a posteriori side all our cognition is piecemeal, and
our thinking discursive; according to comprehension's apriority the
piecemeal must in turn be annulled in the scientific ground-cognition. As
there must accordingly be empirical sciences, which rest upon the a
posteriori and follow their own independent method, so there must
necessarily also be a science of the sciences, namely philosophy, which,
according to its peculiar method, aims at cognizing truth in its apriority.
The concrete carrying-through of philosophy is won only through experience,
and the light that clears up experience streams forth in turn from
philosophy.

\anm{5} The question of the subjective concept's relation to empirical
actuality transforms itself, in the highest instance, into a problem
concerning the relation between \emph{thinking and being}. According to the
subjective idealism, real being perishes in formal thinking; according to
the crass empiricism (sensualism), free thinking goes to ground in
contingent being. Since now each of these theories has shown itself in its
one-sidedness; since subjective thinking cannot produce the fullness of
existence, and objective existence can only be thought, because it is
everywhere subject to the laws of thought, but cannot itself think, since it
lacks subjective freedom: the finite opposition between thinking and being
must be led back to its original ground-unity. In this primal ground, or
original ground, the opposition loses its finitude; for thinking coincides
at every point with being, and both are so far equally infinite. But on the
other side this infinite unity cannot be without difference either. If
thinking and being were in their deepest ground tautological, then the
finite separation, which is precisely what is to be explained from the
ground, would be without any ground. The result is accordingly this, that
the finite, \emph{consecutive} opposition between thinking and being seeks
its sufficient, \emph{primitive} ground in their infinite, absolute
opposition, i.e.\ in that concept in which their separation is as infinite
as their union, so that their outer difference is transformed into an inner,
\emph{living} unity. The absolute Being, which by virtue of its free
thinking produces the actual existence, reveals at the same time the reality
of its own existence in the product of its creative freedom. At this
turning-point the \emph{logical} passes over into the \emph{theological}.
Anselm's ontological proof of God's existence is a logical presentation of
the absolute concept's infinite apriority. The objections that at different
times (Gaunilo---Kant) have been brought against this proof strike not the
substance but the form. For the proof's adequate form can properly be
carried through only in ontology, or the objective logic.

\parag{5}{The Totality of the Concept.}

The complex of all determinations belonging to a concept makes up its inner
totality or \emph{content}; the number of the existing objects in which the
concept finds its expression determines its outer totality or
\emph{extension}. With respect to content and extension, concepts have been
divided in formal logic into \emph{higher} (\textit{conceptus superiores})
and \emph{lower} (\textit{conceptus inferiores}), such that the former,
which have less content but greater extension, comprise the latter under
themselves. Since this division applies especially to the natural relation
between genus (\textit{genus}) and species (\textit{species}), the
superordinate concepts have also generally been called
\emph{genus-concepts}, the subordinate \emph{species-concepts}.

The dialectical relation between the concept's inner and outer totality led
the Schoolmen of the Middle Ages to a metaphysical inquiry concerning the
universality of concepts (\textit{universalia}) and the singleness of things
(\textit{res, ens, entitas, hæcceitas}). The nominalists
(\textit{universalia post rem}) misjudged the concept's essential apriority;
the one-sided realists (\textit{universalia ante rem}) overlooked the
\textit{principium individuationis}. The true realism (\textit{universalia
in re}) is at the same time the highest idealism (\textit{universalia in
mente divina}). The divine thinking, which reflects the absolute concept, is
one with the creative omnipotence, which brings forth that actuality in
which the concept's inner totality is sundered and transformed into an outer
totality of relative concepts, whose existence is posited in and \emph{with}
the existing things.

\anm{1} The actual concept is a thought-whole that finds its expression in a
corresponding sensory whole; it is an inner that comes forward in its outer.
If now the relation between the inner and the outer is to be more closely
determined, these principal sides again step apart from each other as
separate totalities. The outer, which is sensed but not thought, is
indifferent to the inner, which is thought but not sensed; a thing's
empirical or palpable outward can never be derived from its concept: the
intellectual and the sensibly real are so far to be regarded as
incommensurable magnitudes. But on the other side the concept cannot step
out into actuality without presupposing a stuff or material in which it is
realized and comes into view: every outer must answer to \emph{its} inner,
and the object's empirical totality satisfy the concept's demands. If the
surrounding world is regarded with a sensory eye, it is wholly without
concept; if on the contrary it is regarded with the eye of thought, the
concept comes into view in the most various forms. The inadequate relation
between concept and actuality makes itself known not merely in the
normal-concepts being hindered in their full emergence (for example, the
perturbations of the planets; the human organism's want of perfect health),
but shows itself especially in this, that the concept's content, owing to
the outer separation of finite things, falls out of itself into so sporadic
a multiplicity of partial concepts that subjective reflection is not able to
discover their concrete connection. Instead of reducing the partial concepts
to moments in the absolute concept's inner totality, subjective thinking
sets about grouping them according to the various more abstract or more
concrete fundamental forms, and thereby making distinct to itself the
concept's outer totality. The \emph{insight} into a concept's essential
determinations (content) is laid at the ground of a \emph{survey} of the
number of the objects that are subsumed under it (extension). It lies then
in the nature of the case that content and extension stand in inverse
relation.

\anm{2} If the concepts are compared among themselves with respect to
extension, the higher concept cannot be said to contain the lower within
itself, just as little as it exhausts all the intellectual determinations of
the single thing. The higher and the lower concepts then separate from one
another, owing to the objects' outer existence, in such wise that the
former, as the more abstract, are indifferent to the latter; they stand to
one another as wider and narrower schemata, so that the whole comprehending
is reduced to a mere schematizing.

% The Sibbern quotation below CLOSES with '' but has no opening quote anywhere
% — a printer's defect on Danish p. 29, deliberately mirrored from the
% transcription (see its note at that site).
There are cases in which subjective thinking is not able to derive all the
object's special concept-determinations from one fundamental determination,
but must seek different points of departure for its consideration, and
content itself with a combination of several concurrent concepts, each of
which belongs to a different sphere. Thus Prof.\ Sibbern adduces (\S\,42)
that the intellectual penetration into the nature and essence of the
crystallization-forms of a mineral has its own point of support in the
mathematical consideration, whereas its investigation in respect of lustre,
color, and so on has quite another.'' Insofar as a fundamental concept's
modifications arise from outer concurrences and empirically determined
circumstances, they can be overlooked as unessential and indifferent; but
what is thus in one respect unessential is, from another point of view,
under the presupposition of another fundamental concept, to be regarded as
essential, and the various regards and combinations that the empirical
pursuit of thought requires do not concern logic as such.

If on the contrary the concepts' own relations are developed, it appears
that the higher concepts always have their specific content-determinations
in the lower. This is already expressed when the superordinate concepts are
determined as genus-concepts (γένη), the subordinate as species-concepts
(εἴδη); for here it holds that every genus must contain its species. Life,
regarded as genus-concept, requires not merely a plant-world and an
animal-world, but also a human world for its complete expression. Organic
nature presents this in such wise that the lower potencies of life, which
win a sporadic existence in separate groups of individuals, whose genera and
species are in turn relative concept-determinations, make up the basis upon
which a higher potency of life develops itself. Thus man has the vegetative
and animal life as basis for the rational. The partial appearances in which
the concept of life is sundered and realized on the lower stages are
gathered up and taken in as moments in the higher totality of human life.
The proposition that man shall comprehend his life rests in the
presupposition that the concept of life shall be realized in man.

\anm{3} When it is said that the concept, as such, is \emph{the universal},
this proposition acquires an entirely different meaning according as
subjective thinking comes to cognition of the universal by an
\emph{abstraction} or by a \emph{reduction}. If the universal is held fast
by abstraction, then the highest concept is reduced to an empty schema, an
intellectual side, a common mark, which serves the reflection that orients
itself in the outer multiplicity of things as a mere \textit{ratio
cognoscendi}. But when the universal is cognized in its true universality,
it is not merely a side of everything, but the essential allness itself, the
universal that penetrates everything, actualizes itself in everything, so
that all single thoughts, all existing things, must let themselves be
reduced to moments in its concrete content. In this sense the general
concept is apprehended as the absolute concept, as the ground-unity that
contains the \textit{ratio essendi} of all separate concepts.

When it is further said that the concept is the \emph{separate}, it is
thereby expressed that multiplicity, difference and opposition belong to its
intellectual essentiality, that the specific peculiarities and changing
modifications (διαφορὰ, ἴδιον, συμβεβηκὸς) whereby things separate
themselves from one another fall under thinking's point of view. Yet in this
respect too it makes an essential difference whether the subject comes to
cognition of separateness by mere \emph{comparison} or by an immanent
\emph{deduction}. For the outer comparison the separate concepts fall apart
from one another, so that their existence must be sought in the sporadic
groups of objects, and \emph{combination} unites them according to various
regards. Universality is here again the empty schema. Deduction, on the
contrary, presupposes the fundamental concept's concrete universality, from
which all separate concepts and concept-moments are derived. According to
outer comparison there is an innumerable host of partial concepts; according
to deduction the general concept is itself sundered into its separate
determinations, and thereby completed in a \emph{logical concept-system}.
The more comprehensive the deduction is, the more it can make mere
comparison and combination superfluous, the higher lies the concept's
sphere, and the deeper is the insight.

When it is finally said that the concept itself is the \emph{single}, a
distinction must be made between the \emph{sensory} and the
\emph{intellectual}, between the \emph{immediate} and the \emph{reflected}
singleness. That the sensory singleness (that is to say here, the object in
its merely empirical wholeness) does not as such concern the concept has
already been shown in the foregoing, and becomes still more evident when we
determine more closely the relation between the species-concept and its
individuals. In relation to the genus, the species-concept is posited as a
\emph{separate} concept, in which the genus-concept's general content is
specified; in relation to its individuals, the species-concept must on the
contrary be developed into a single concept specified within itself, which
as such is indifferent to the sensory singleness. The distinction Porphyry
(in his introduction to the Aristotelian categories) has made between γένος,
εἶδος, ἰδιον on the one side and διαφορὰ, συμβεβηκὸς on the other --- that
the former, namely, are used of a thing's ``Qvidität'' (ἐν τῷ τί ἐστι
κατεγορεῖται), the latter of its ``Qvalitet'' (ἐν τῷ ὁποῖόν τί ἐστι κατ.)
--- is a merely empirical distinction, which falls away when comprehending
thinking is carried through. The peculiar (τὸ ἴδιον) is just as much a
qualitative as a ``quidditative'' determination, and the intellectual
difference (ἡ διαφορὰ) is itself the moment whereby the transition is made
from the genus-concept to its species-concept. The specification therefore
cannot stop before the species-concept is absolutely individualized; but
such a \textit{species infima} is a single concept that maintains itself in
the changing exemplars, insofar as these are apprehended merely from the
sensory side. --- Since now the genus-concept is always foundation for the
species-concept, and this in turn is specified until the rational
individuation is completed, the universal (the genus) and the separate (the
species) merge into the concept's intellectual singleness (the logical
individual).

The universal, the separate and the single concept are then, in consequence,
not to be regarded as three sorts of concepts, since they much rather
present the concept reflected in absolute thinking from three principal
sides; but just as little is their distinction an empty, formal one; for in
finite existence they fall apart from one another under innumerable
modifications and relations, and shape themselves into a world of real
concepts.

\anm{4} Since the actual world presents itself at once both as an infinite
complex of empirical objects, in which the concept is lost, and at the same
time as an intellectual world, i.e.\ as a world of concepts, in which the
sensory things vanish as indifferent magnitudes: there arises the
metaphysical problem whether the reality of concepts is primarily to be
sought in things, or the reality of things in concepts. While Plato declared
himself especially for the latter, Aristotle on the contrary approached the
former of these views, and Porphyry (Ειςαγωγη: first section) puts the
% The defective Greek title (no breathing, no accent, final sigma set
% medially) is as the Danish page prints it; copied verbatim from the
% transcription, per its note.
problem in the following terms: αὐτίκα περὶ γενων τε καὶ εἰδῶν, τὸ μὲν εἴτε
ὑφέστηκεν, εἴτε καὶ ἐν μόναις ψιλαῖς ἐπινοίαις κεῖται, εἴτε καὶ ὑφεστηκότα
σώματά ἐστιν, ἢ ἀσώματα, καὶ πότερον χωριστὰ, ἢ ἐν τοῖς αἰσθητοῖς, καὶ περὶ
ταύτα ὑφεστῶτα, παραιτήσομαι λέγειν, βαθυτάτης οὔσης τῆς τοιαύτης
πραγματείας. Through the whole Middle Ages philosophy proper was occupied
almost exclusively with the solution of this problem. The crass nominalism
confused the concept's existence with the sensory singleness, negated the
reality of concepts, and posited them as empty schemata (\textit{nomina,
flatus vocis}). The one-sided realism went to the opposite extreme, annulled
the intellectual singleness, in that it wanted to negate the independence of
the sensory things, and placed reality in the general fundamental forms and
unchangeable types. The one-sidedness of both theories was demonstrated by
Abelard. The fundamental thought in Abelard's demonstration is this, that
neither party is able to give any true definition of the essence of a thing.
The nominalists could not define, because they held fast the single without
acknowledging the universal. For since definition rests upon an essential
difference between \textit{genus} and \textit{species}, every nominalist
definition that denies this difference must be wholly indifferent. If the
concept \textit{animal} is a mere \textit{nomen}, then the concept
\textit{animal rationale} is also only a \textit{nomen}, and the difference
between beast and man consequently purely nominal. Just as little could the
one-sided realists arrive at any true definition; for since they overlooked
singleness, they had to confine themselves in every proposition to
predicating the universal of the universal, whereby the concept's
determinateness was wholly lost. If the man in Socrates, for example, were
also man in all other individuals, then it would follow that the whole of
mankind must be sick every time Socrates was sick. These extremes the
medieval philosophers, with all their reasoning acumen, could not reconcile.
As realism developed itself to an acknowledgment of individuation, it had to
seek the concepts' existence in things, and thus came to set up the theory
of \textit{universalia in re}. But when the existence of concepts and of
things immediately coincides, thinking must again begin with singleness, and
abstract the concept from the thing. The concept in its intellectual purity
thus has validity only for abstracting and reflecting thought. The doctrine
of \textit{universalia in re} led in this way to the doctrine of
\textit{universalia in mente nostra}, whereby the developed realism passed
over into a higher nominalism. The true yield of the Schoolmen's inquiries
is the recognition that the concept's all-sided existence requires an
intellectual as well as a sensuous medium. With intellectual freedom the
concept exists only in comprehending thought (\textit{in mente}); with
immediate actuality it first comes forward in its sensory material
(\textit{in re}). But if its reflected being \textit{in mente} is now to let
itself be reconciled with its immediate existence \textit{in re}, then both
these extremes must seek their original ground in the divine thinking, which
is one with the creative productivity. The theory of the empirical realism,
that the concept exists \textit{in re}, must then be raised to the truth
that it is actualized by the comprehending omnipotence, which is the free
cause of things; and the nominalist proposition, that \textit{universalia}
exist merely \textit{in mente nostra}, subordinated to that proposition
which is idealism's highest expression: that they exist \textit{in mente
divina}. Thus the subjective logic points once more to the objective.

\secrule

\capitel{Second Chapter.}{The Presentation of the Subjective Concept: the Logical Judgment.}
\parag{6}{The Logical Form of the Judgment.}

Language is the intellectual medium in which the thinking subject possesses
its own concept-content and that of the surrounding world; hence the
concept's rational purity and freedom can first be presented in words and
sentences. Since no sentence may be meaningless, the concept is already
discernible in every intelligible utterance; but if the meaning, as the
merely represented concept, is to be expressed in the concept's own form,
then the sentence itself must be transformed into a logical judgment. As the
judgment's \emph{single concept} (the subject) is subsumed by the copula
``is'' under its \emph{general concept} (the predicate),
\emph{separateness}, or the determinate concept, comes into view in the
difference and unity of the two. The unity of subject and predicate is an
\emph{identity}, which as such contains its difference; but since this
difference is in turn different from the identity, \emph{contradiction}
enters.

The development of the mutual relation and logical significance of the
principle of identity and the principle of contradiction affords one of the
most essential points of dispute between the formal and the speculative
logic.

\anm{1} Every object that answers to its concept can at the same time be
said to \emph{present} the concept. The simplest way in which a thinking
subject can show that it has some concept of the thing will thus be to
point out the objects in which the concept is to be sought. This mimic
presentation is, however, highly imperfect, and proves that the subject has
not got beyond ``the general representation'' and ``the clear and distinct
intuition.'' Hence it is also a criterion of merely representational
thinking that it more easily produces an example than an adequate
definition. --- The concept is appropriated with greater freedom when it is
possessed as guiding norm for an independent construction. The so-called
``construction-concepts'' (jfr.\ Sibberns Logik Pg.\ 137--163) may properly
be said to lie on the boundary between theory and practice. The presentation
of the concept here coincides immediately with its actualization. But the
mechanical productions, considered in and for themselves, are only ambiguous
testimonies of subjective insight, since bees, artisans and automata can up
to a certain point present the concept in the same style. Of thoroughgoing
significance, on the contrary, is the distinction between the artistic and
the scientific, the plastic and the logical presentation of the concept's
content. The plastic presentation keeps closest to \textit{universalia in
re}, the logical closest to \textit{universalia in mente}. The artist, whose
thoughts are intellectual visions, expresses the concept \textit{in
concreto} by presenting its image. That the concept really is in movement
during artistic production is evident from the fact that art's ideals can be
developed in the theory of art, so that aesthetics constitutes an integral
member in the system of philosophy. But the plastic presentation of the
concept's content shows nonetheless its specific difference from the logical
in this, that the artist's genial productivity by no means rises in the same
degree as his reflecting cognition of the concept of art develops. The more
sensuous art's material is, the more indifferently can it stand toward
science; the more intellectual, on the contrary, the artistic medium shows
itself, the more distinctly will the interaction between the logical and the
plastic presentation come into view. Hence poetry enters into the most
inward relation to philosophy; for poetry is the art of the word, and the
word is the medium not only of intuition and representation, but also of
thought and the concept.

\anm{2} In Plato's dialogue ``Cratylus'' the problem is treated whether the
word, as such, is originally to be regarded as a natural copy of the thing
(μίμημα), or perhaps rather as an arbitrary sign (θέσει); this is, in other
words, the question whether language may be assumed to be a mere product of
nature, or a means of communication chosen by common agreement. This inquiry
concerning the origin of language concerns logic only insofar as the
relation between the thinking subject's receptivity and spontaneity, between
word and thought, representation and concept, hereby receives a closer
determination. The Heraclitean view, that the sound of the word is
originally a sensory imitation of the thing, corresponds to the sensualist
theory of the \textit{tabula rasa}. The soul, whose self-activity is lost in
feeling and receptivity, is regarded as an instrument whose tones are
produced by the impressions of the surrounding world, and the word is
confused with the immediate sound of nature, which the beast has in common
with man. The Democritean and sophistic theory of the understanding, which
denies language's nature-side, is akin to the formal idealism, and moves in
a tautological circle. The arbitrary choice of the word presupposes a
reflection, but reflection in turn presupposes language, for all thinking is
a speaking with oneself. In order to reach an agreement concerning an
intellectual means of communication, the individuals would already have to
be in possession of a means of communicating their thoughts to one another.
The origin of language points to an immediate unity of sensation and
thinking, a living interaction between the subjective and the objective
reason, and can therefore, in logical generality, be regarded as a spiritual
product of nature, which in the development of culture is more and more
penetrated by free reflection. While Plato especially brings out language's
nature-side, Aristotle on the contrary urges the culture-side, and regards
every noun as a free designation (φωνὴ σημαντικὴ κατὰ συνθήκην) of the
representation in which the soul holds fast the image of a determinate
object. The single principal words of language can therefore be apprehended
from a double point of view: partly they are to be seen as involuntary signs
(τῶν ἐν τῇ ψυχῇ παθημάτων σύμβολα), to which the imagination has bound the
images of things; partly they are reflection's own thought-symbols, in which
it unceasingly hovers between representation and concept. It is often the
case that a concept does not break through, and come to clear development,
because it lacks its corresponding word --- just as, conversely, the most
primitive and pregnant expressions gradually fade and lose their freshness
when the thought-content whose bearers they originally were is diluted and
trivialized. (Philosophical root-words, categories --- phrases.)

\anm{3} Every word carries a representation with it; but the word standing
alone gives as yet no meaning. Meaning first arises when the word that
immediately designates the thing's existence (the noun) is joined with
another (the verb), which expresses its being and the utterance of its
existence. Such a combination of words is a \emph{sentence}, which again can
be considered from the grammatical and from the logical side. From the
logical point of view, the same relation obtains between the mere sentence
and the judgment proper as between meaning and concept. As the subject in
the sentence may be an accidental representation, so also the predicate may
express a merely passing determination without the stamp of universal
validity. Sentences such as: the boy plays, the snow melts, \&c., cannot
therefore without further ado be reckoned among the logical judgments. The
formal mark whereby the judgment distinguishes itself from the mere sentence
is not primarily to be sought in the subject, but rather in the copula and
the predicate. If the thing is held fast in its immediate utterance or
appearance, the copula can be contained in language's many different active,
passive and neuter verbs (cryptic sentences). But if predication, as an act
of thinking, is to come forward with logical universality, then it is at
once evident that everything that can be affirmed or denied of a given
object must in the last instance refer to being and not-being. Since now the
sentence: \textit{A} is, can be opposed both by: \textit{A} \emph{was}, and
by: \textit{A} is not, it follows that the copula ``is'' must distinguish
itself from the verb in this, that tenses and moods lose their significance
in the logical present of the position. As for the predicate, it belongs to
its essence to express the universal, as such. Hence such utterances as: the
rose is red, the negro is black, \&c., are likewise not to be regarded
outright as judgments proper, although from the side of form there is
nothing to object. Everything here depends upon the significance that is
ascribed to the grammatical predicates red, black, \&c.; if these are
apprehended merely as immediate sense-determinations, then the utterances
named are only sentences; if on the contrary they are \emph{thought} as
universal determinations in opposition to the subjects' singleness, then we
have the judgments proper. That the predicate presents the concept's
universality must not be understood as though it could express the concept
only \textit{in abstracto}, while the subject gave it \textit{in concreto}.
In such judgments as: ``Man is rational,'' ``The beast is living,'' the
predicate is certainly more abstract than the subject, but precisely for
that reason these judgments are not to be regarded as the concept's adequate
expression either. In adequate judgments, subject and predicate have the
same content (for example, God is the absolute Spirit). The reason why the
subject is nevertheless apprehended as single concept, the predicate as
general concept, is this, that the subject refers proximately to
representation, while the predicate refers to thought. Insofar as the
represented concept is at the same time possessed as a concept set out in
existential immediacy, the subject expresses its \textit{universale in re},
which the predicate in turn annuls in reflection, and posits as
\textit{universale in mente} (Stilpo's cabbage). With respect to the
concrete subject, the concept's universal determinations are concentrated
into immediate singleness; with respect to the developed predicate, the
single concept's content is reflected through into pure universality. Both
principal sides of the concept are thus each for itself posited as the whole
concept, and yet these two are, by means of the copula ``is,'' united in one
positing, which expresses the mutual and thereby separate determinateness of
both. This separateness comes into view in every judgment; for when it is
said that the rose is red, there is just as much talk of the red rose as of
the rose-red.

\anm{4} In every judgment the predicate is to declare \emph{what} the
subject is. Insofar as the judgment is not entirely adequate, the subject
may well contain more determinations than the predicate; but this difference
does not annul the likeness between subject and predicate, insofar as the
determinations falling outside the predicate must in all such cases be
assumed to fall outside the judgment as well. To hold fast the unity of
subject and predicate, formal logic has set up the so-called principle of
identity: \textit{A} \emph{is} \textit{A}, \emph{every thing is what it
is.} But on the other side the predicate must nevertheless add a new
determination to the subject, or at the least render the determination in a
new form; he who asks what \textit{A} is cannot be satisfied with the answer
that \textit{A} is \textit{A}, for that he knew already beforehand. By
urging the formal principle of identity, the Megarian Stilpo contested all
judgments proper. One could not, in his opinion, say: ``The horse runs;''
for ``horse'' and ``to run'' are not identical, since the same predicate is
ascribed also to dogs and lions. Nor could one think the sentence: ``Man is
good,'' for ``good'' and ``man'' are not the same, since the predicate
``good'' is predicated also of bread and medicine (Poul Møllers efterladte
Skrifter, 2det Bind Pg.\ 383). Formal logic therefore also concedes that
every thing must be determined by its relations to other things, that
identity as such contains difference, that he who directly declares what the
thing is thereby also indirectly declares what it is \emph{not}. But as the
subject's identity with itself passes over, by means of the predicate, into
difference, thinking must at the same time negate every difference that will
not merge into the identity, and thereby hold fast the \emph{agreement}
between subject and predicate. To this end logic sets up its formal
\emph{principle of contradiction}: ``\textit{A} is not that which \textit{A}
is not;'' or: ``No thing is that which it is not, and is not that which it
is;'' or: ``No thing can have predicates that contradict the thing;'' or:
``All pursuit of thought shall be consistent;'' or: ``The contradictory is
at the same time the unthinkable.'' What the principle of identity is in the
positive form, the principle of contradiction is in the negative. The formal
advance consists in this, that the predicate's agreement with the subject,
which according to the fundamental proposition of identity is assumed as
immediately being, becomes reflected through by the application of the
principle of contradiction, and that subjective thinking has found in this
reflection its proper touchstone.

\anm{5} As the contradiction is the logical sting that wakes thinking from
its slumber, and stimulates it to independent activity, so all true thinking
also aims at overcoming and annulling the contradiction, in order to come
again to rest and thorough satisfaction. Could thinking find rest in the
unresolved contradiction, then black would be white, the circle
four-cornered, truth a lie, the good evil, and the highest philosophy a
sophistical deception. When the polemic between the defenders of the formal
and of the speculative logic has at times taken on the appearance as though
the former merely wished to exclude the contradiction, the latter to adopt
it, this tension rests only upon a misunderstanding, which a closer
consideration must make disappear. In formal respect, each of the contending
parties is compelled to concede to the contradiction the same significance.
The defenders of formal logic must necessarily take the contradiction up
into their thinking, immerse themselves in it, in order to be able to
contest it. For it is not with the contradiction as with an outward enemy,
whom the well-informed can avoid and thereby escape the danger; the
contradiction pursues the thinker at every step, so that the true
development of thought, as regards its formal consistency, consists
precisely in contradicting the contradiction everywhere, or, what is the
same, that all thorough thinking goes through sheer contradictions. On the
other side, the speculative logic cannot defend the contradiction without
also attacking it. For it wants precisely to prove that formal logic, by
setting up its principle of contradiction and denying the contradiction its
logical right, makes itself guilty of an indefensible contradiction.
Regarded from the purely formal side, the inquiry is thus rather unfruitful;
but the main thing is that the difference between the logic of the
understanding and the speculative logic of reason be clearly and
determinately apprehended. This real difference will come more closely into
view when, from the double standpoint of the doctrine of thought, we
consider: a) the concept of the contradiction, b) its relation to the
thinking consciousness, and c) its scientific solution.

\begin{itemize}
\item[a)] When the concepts are placed under the point of view of finitude,
  they can be separated by determinate definitions, and the rule holds that
  they must not be confounded; if on the contrary they are regarded in the
  light of infinity, then the one concept must be completed by the other,
  since every isolated concept-determination is defective and one-sided.
  This double way of seeing leads us to a determinate distinction between
  the \emph{superficial} and the \emph{thorough} contradiction. To confuse
  the like with the unlike, unity with multiplicity, or the whole with the
  parts, is to make oneself guilty of a superficial contradiction; but he
  who proves that likeness cannot be without unlikeness, that the concept of
  unity contains that of multiplicity, thereby points out the thorough
  contradiction. The difficulty of breaking down representation's finite and
  immediate barriers, without at the same time volatilizing the concepts'
  essential difference, stood clear before the Greek philosophers, who
  recognized how easily a dialectical ἀδολεσχία passes over into φλυαρία.
\item[b)] It lies then in the nature of the case that formal logic, which
  apprehends the concepts in their finite relations, must polemicize against
  the contradiction and regard it as a fault in subjective thinking, while
  the speculative logic on the contrary refers the contradiction to
  existence's own ambiguity, and sets a considerable strength of thought in
  drawing it forth. As formal logic makes only a \emph{critical} application
  of the contradiction, it proceeds from the presupposition that what
  contains a contradiction cannot exist. But that this way of seeing is
  insufficient is already evident from the fact that the thinking of the
  understanding, under the most various turns, has had to admit to itself
  that the very highest truths (creation, freedom) appear to it
  contradictory and incomprehensible. The speculative logic, which makes an
  \emph{architectonic} use of the contradiction, can on the contrary employ
  it as a guide from the one sphere of existence to the other. In Plato's
  philosophical wonder, in the mystics' spirited paradoxes, in the Kantian
  antinomies, we everywhere sense the contradiction's dialectical
  pulse-beat. Without a philosophical problem there is no philosophical
  productivity; but to pose the true problem is, in other words, to
  apprehend the thorough contradiction.
\item[c)] The fundamental proposition of formal logic: \textit{Diversi
  respectus tollunt contradictionem}, must in the speculative logic be
  converted into: \textit{Diversi respectus faciunt contradictionem}. To
  solve the contradiction by the help of \textit{diversi respectus} is the
  same as to let the different regards fall away from one another and lose
  themselves in a reasoning of grounds; to reduce \textit{diversi respectus}
  to their contradiction, in order to let them vanish in the deeper
  fundamental view, is to seek the absolute ground. Formal logic, which
  merely teaches us to think \emph{about} the objects, takes these as they
  are given, and ascribes to things an existence closed in itself and
  indifferent to our subjective movements of thought. Under this
  presupposition it has its complete validity that the matter in question
  must be posited as τὸ αὐτό and regarded κατὰ τὸ αὐτό, πρὸς τὸ αὐτό,
  ἐν τῷ αὐτῷ χρόνῳ. But if it is already a questionable matter that the last
  of these demands properly cannot be fulfilled at all, since every Now
  unceasingly vanishes, then it is a still more essential defect that the
  subject, which sets the one regard up beside the other in order to hold
  fast an incessant \textit{idem per idem}, must remain standing at
  consideration and description, without reaching thinking proper. The
  speculative logic, which thinks through the genetic development of things
  and concepts, must therefore unceasingly set its paradox against the
  vulgar tautology. It is not denied that something, insofar as it is
  posited as one, is regarded from another point of view than when it is
  posited as several; or that what is determined as ground is determined
  according to another regard than when the same is determined as
  consequent. But what is most determinately denied is that these regards
  should fall indifferently apart from one another. There is \emph{both}
  identity and difference; but both proceed from one and the same ground;
  the mere identity without difference is just as contradictory as the mere
  difference without identity; there is identity \emph{because} there is
  difference; the determination that hits the identity hits also the
  difference: the identity is determined \emph{as} difference. There is in
  the judgment \emph{both} a subject \emph{and} a predicate; but the subject
  is what the predicate is, precisely \emph{because} subject and predicate
  are \emph{not} the same. In the unceasing advance of thought and
  existence, or in the living movement, lies then the true solution of the
  contradiction.
\end{itemize}

\parag{7}{The Logical Content of the Judgment.}

Since the judgment has the significance of being the concept's logical
presentation, it follows that its formal development cannot be separated
from the concept's content. Insofar as the single judgment, considered in
and for itself, expresses only a partial thought-result, it presents the
concept only from a certain side, and the many isolated judgments can
therefore be classified according to very different regards. Formal logic
has divided the judgments: a) according to \emph{quality} into: affirmative,
negative and infinite; b) according to \emph{quantity} into: singular,
particular and general; c) according to \emph{relation} into: categorical,
hypothetical and disjunctive; d) according to \emph{modality} into:
assertoric, problematic and apodictic judgments. But these divisions have
logical worth only insofar as they also show the gradual rising of insight
and of the power of judgment.

By the help of the logical judgments, the speculative logic develops the
concept's own inner and outer sunderings or \emph{diremptions}.

\anm{1} That the judgment is more than a merely subjective act of thinking,
by which reflection satisfies itself without regard to objective being, is
already expressed by the copula. Nevertheless, the separation and union of
subject and predicate, as it takes place in the judgment, is not a movement
that existence itself realizes: the immediately existing concept passes no
judgment upon itself; the act of judging goes on in the thinking
consciousness. Only the thought concept is as such the judgment's content;
but this is in turn identical with the concept that is, or exists. Hence
every judgment's logical dignity must be determined as well by the power of
judgment it requires in the subject as by the objective concept-content it
presents. It is then self-evident that the judgment's logical content must
indeed be distinguished from its matter or empirical content. In such
judgments as: Caius is a man, Pasop is a dog, the matter is indeed
different, but the logical content still the same.

\anm{2} The qualitative judgment posits the thing as subject and the
property as its predicate: the rose is red: the single is the universal.
This identity of the single and the universal is as yet quite abstract, and
therefore passes over into an equally abstract difference. The red rose is
immediately determined by its color; but partly the rose itself can be
thought to exist with another color, and still be a rose; partly the color
can be met with in other objects, and still be red: the cheek is red, the
wine is red, and so forth. In the qualitative judgments a contradiction
therefore comes into view, which has its ground as much in existence as in
thinking. As regards existence, the one thing is always determined in
relation to the other; what the one object is in and for itself, it is in
relation to the other, and it is for that reason continually subject to
change:
``At this moment the water is solid, in the next moment fluid,'' ``The leaf
that is now green will in some time be yellow.'' The existing subject
unceasingly changes predicates, and the transience of things shows itself as
an unbroken hovering between \emph{being} and \emph{not-being}. In
consequence, judging thought must likewise move between \emph{affirmation}
and \emph{negation}, and separate the judgments of quality into
\emph{positive} and \emph{negative}. This separation, however, when more
closely regarded, is no proper division; for the positive judgment shows
itself in its abstract positivity just as one-sided and untrue as the
abstractly negative. The single (the concrete thing with all its properties)
cannot be immediately identical with the abstractly universal (the single
property). The positive judgment can therefore be resolved a) into two
% The Greek "β)" standing where "b)" is wanted is the Danish page's own wrong
% sort (see the transcription's note at p. 48); mirrored as printed.
tautological sentences: the single is single, the universal is universal,
and β) into two superficial contradictions: the single is a universal, and
the universal is a single. Just as little can the merely negative judgment
be isolated from the positive. When it is said: the rose is red, it is still
presupposed that it has a color. As the negative judgment negates a
determinate species-concept, this always happens under the presupposition
that the genus-concept lying at the ground can be determined in another
species. If on the contrary the single is held so isolated that the
genus-concept is negated together with the species-concept, then the
judgment itself falls outside the concept's sphere: the so-called
\emph{infinite} judgments (conscience is no elephant; reason is not
four-cornered) are only thoughtless sentences. That the positive and the
negative judgment must reflect each other is evident also from the fact that
only a very slight degree of the power of judgment is needed in order to
utter them each by itself; whereas the proper insight into the concept's
essence rests precisely on this, that he who directly passes the positive
judgment has the negative \textit{in mente}, or conversely is conscious that
the negative judgment holds only in virtue of the positive.

From the standpoint of the speculative logic, the judgment of quality shows
itself as the existing concept's own diremption (Hegel: Das Urtheil des
Daseyns). The existing things pass away; only the concepts in their
universality endure: \emph{the single, which is a universal,} is dissolved
and vanishes in the stream of universality. But \emph{the universal is also
a single}: the concept expresses its existence in the change of things;
hence single existence is reproduced just as unceasingly as it vanishes.
--- This diremption of the concept casts a new light upon the strife of the
nominalists and the realists. The crass nominalists apprehended the
qualitative judgment according to its immediate content, and thereby came to
posit a multiplicity of subjects (\textit{res}) without predicates; the
one-sided realists, who kept to the judgment's abstract form, dissolved
existence into an aggregate of predicates (fluttering \textit{universalia})
without determinate subjects (\textit{res}).

\anm{3} The qualitative judgment declares that the universal is to be sought
in the single, that the predicate ``\emph{inheres in}'' its subject, while
on the other side it is recognized that the subject is to be subsumed under
its predicate. Schematizing thought therefore undertakes to test the
predicate's extension by determining the number of the subjects. This test
is instituted in the judgments of \emph{quantity}: ``\emph{This} single,
\emph{some} singles, \emph{all} singles are comprised under the
universal.'' But such a counting-up has only a faint shadow of comprehension
to show; reflection loses its way in an outward comparison between the
universal and the single. If the quantitative judgment is to carry with it
any true advance in the cognition of the concept, then the comparison must
pass over into a subsuming reflection, whereby the judgment itself is to be
characterized as a \emph{comprehensive} or \emph{comprising} judgment.
According to the qualitative judgment, the subject is immediately set out as
the substrate (ὑποκείμενον) to which the predicate attaches; the predicate,
which is chosen immediately, must therefore run through the reflection of
affirmation and negation. But hereby it takes on a stamp of essential
self-validity, which is presupposed in the comprehensive judgment, whose
subjects are comprised --- singly, partially or generally --- under their
predicate. Hence Hegel calls attention to the fact that the predicates
adduced in examples of the comprehensive judgment (``Das Urtheil der
Reflexion'') must be of another kind than those suited to quality. Such
sentences as: man is mortal, things are useful, and so on, can serve as
examples. The power of judgment that the comprehensive judgment requires can
be determined as merely reflecting: through the subjects' separate groups
thought glimpses the predicate's essential universality. The concept itself
hereby comes into view as the \emph{relation} of the existing things. But
this relation has a double side, since it is both changeable, as the things'
phenomena, and at the same time unchangeable, as the things' essence. If the
relation is apprehended from its unchangeable side, the concept is
determined as the things' law, from which there is no exception; if on the
contrary the relation is regarded from the point of view of changeableness,
the concept is abstracted as existence's rule, from which there are now
more, now fewer exceptions. In the merely comprehensive judgment the latter
way of considering is the predominant one: universality is determined only
as allness, and the general is not yet posited as the true universal
(empiricism).

\anm{4} The division of the judgments according to their \emph{relation} has
the logical significance that the general (the collective universality) is
developed into true universality, whereby the relation between the
genus-concept and its species-concepts is essentially determined. The
categorical judgment's predicate expresses the subject's inner nature and
essence. The comprehensive sentence: All speaking animals are men, first
finds its true form in the categorical utterance: The speaking animal is, as
such, man; whereby thinking presses through the empirical mark into the
concept's essence, and finds the thing's category. If the categorical
judgment's content is developed further, herein lies that the thing cannot
express its immediately determinate concept (its species) without answering
to its category (i.e.\ its essential general concept, the genus). This
relation is expressly brought out in the \emph{hypothetical} judgment's
different forms: If \textit{A} is, then \textit{B} is; if \textit{A} is
\textit{B}, then \textit{C} must be \textit{D}, and so on. In this
development, experience passes over into insight and comprehension proper.
The fact vanishes in its contingency, for the question is not primarily
whether \textit{A} or \textit{B} actually exists; but thought's law is
raised to existence's norm; for it is expressly recognized that an
\textit{A} that cannot be \emph{thought} without \textit{B} cannot
\emph{be} without \textit{B} either. --- The real and the intellectual
connection are posited in their essential unity; but this identity, when
more closely regarded, again leaves an essential contradiction behind.
Although \textit{B} is posited in and with \textit{A}, \textit{A} and
\textit{B} are still by no means the same. According to pure thinking they
stand to each other as ground and consequent, cause and effect, but this
relation is as yet wholly formal; according to their existence they fall
immediately apart from each other, and stand as \emph{the condition} to
\emph{the conditioned}. Reflection, which still hovers between the fact's
contingency and the thought-law's necessity, seeks an unconditioned as
foundation for all mutual conditions. This foundation thinking first wins in
the \emph{disjunctive} judgment: \textit{A} is either \textit{B} or
\textit{C}, and so on. The subject \textit{A} now expresses the
genus-concept in its concrete universality, while the genus's determinate
species are separately posited in the predicate: ``Man is (according to sex)
either man or woman.'' The essential relation between the genus and its
species thus finds its logically all-sided expression in the disjunctive
judgment. As the species-concepts taken together are subordinated to the
genus-concept, they show themselves at the same time as mutually
coordinated. The disjunction consists in this, that the species-concepts
exclude one another, because they, each for itself, are the genus's
contentful appearances; they are thus \emph{contrary} \emph{because} they
are \emph{contradictory}, and contradictory because they are contrary.

That the \emph{division} of the judgment, when apprehended from a
speculative standpoint, is a concept-sundering or \emph{primal division}
(\textit{Ur=Deling}), is according to relation wholly evident. In order to
% The Danish page's letterspaced run takes in "har Schelling" as well as "det
% Absolute" — a compositor's lapse logged in the transcription. The emphasis
% is carried here on "the absolute" alone, since the lapse cannot be
% reproduced across the English word order.
make \emph{the absolute} comprehensible, Schelling (``Philosophie und
Religion,'' Pag.\ 11---13) has therefore also availed himself of relation's
forms. ``The first form in which absoluteness can be posited is the
categorical.'' But since the absolute, as such, is \emph{neither} infinite
\emph{nor} finite, \emph{neither} ideal \emph{nor} real, \emph{neither}
subjective \emph{nor} objective, the productive intuition must step in and
fill the negative reflection's emptiness. ``The second form for the
absolute's appearance in reflection is the hypothetical.'' When a subject
and an object are, then the absolute is the identical essence of both.
Subject and object \emph{condition} each other, but the absolute is their
\emph{unconditioned} foundation, their positive unity regarded \textit{sub
specie æternitatis}. ``The third form in which reflection likes to express
the absolute, and which is especially known through \emph{Spinoza}, is the
disjunctive.'' There is only one substance, but this can be regarded either
as purely ideal or as immediately real. Just because the absolute substance,
according to the disjunctive judgment, is either wholly subjective or wholly
objective, therefore it is, according to the categorical form of expression,
neither the one nor the other.

\anm{5} The comprehensive judgment presents the concept's content in
\emph{empirical} sunderings; the disjunctive reduces the empirical content
to a moment in the concept's \emph{logical} sunderings. The intellectual
world and the real world must accordingly be measured by each other: the
thought concept must have its existence, and every existing thing must in
turn answer to its concept. This relation, totalized in itself, is now
further illuminated as we consider the division of the judgments according
to their \emph{modality}. The \emph{assertoric} judgment expresses the
obtaining relation between thing and concept in a wholly immediate way:
``The house is well built; the plan is excellent; the treatise is
scientific.'' But what in one respect answers to the concept can in another
respect deviate from it. Against the simple affirmation stands the
straightforward denial, and the judgment becomes \emph{problematic}. The
problem thus developed can now neither be solved by a merely empirical
demonstration, since all actuality must be valued according to the concept,
nor by a merely formal operation of thought, since the merely thought
concept to which no actuality corresponds is just as untrue as a merely
palpable actuality that contradicts its concept. The insight that actuality
has its truth in the concept, or what is the same, that the concept
completed in its own inner consistency has found its corresponding outer, is
presented in the \emph{apodictic} judgment.

As the difference between the thought concept (\textit{universale in
mente}) and the existing concept (\textit{universale in re}) comes into view
in modality, so modality in turn can be apprehended from a subjective and an
objective standpoint. --- It is characteristic of the Kantian philosophy
that it regards the categories of modality as ``the postulates of empirical
thinking'' (jfr.\ Kritik der reinen Vernunft, Pag.\ 217---223), whose
determinations do not concern the object at all, but fall wholly within the
thinking subject's own reflection. What agrees with experience's formal
conditions is \emph{possible}; what is in connection with its material
conditions is \emph{actual}; what finally is determined according to the
essential unity of both kinds of conditions is \emph{necessary}. By the help
of the problematic, assertoric and apodictic judgments Kant measured the
certainty of human cognition, and determined the logical relation between
\emph{opining}, \emph{faith} and \emph{knowledge}. What is true in the
Kantian theory is that the concept can, without regard to its material
existence, have a purely logical existence in the subjective consciousness
(\textit{universale in mente}); but the system's nominalist limitation shows
itself in this, that the absolute concept, whose infinite validity, grounded
in the divine thinking (\textit{universale in mente divina}), is apodictic,
is confused with the relative concepts, which have an abstract being
\textit{in mente nostra}, and whose real existence must therefore always be
regarded as problematic, until experience brings a decisive certainty. ---
Hegel has therefore rightly transformed the division of the
modality-judgments into the concept's own primal division (Das Urtheil des
Begriffs). The fundamental thought herein is that a concept that exists only
in our own consciousness does not deserve the name of true concept. It is
then either a mere representation, an imaginary magnitude, of which the
imagination can produce as many as you please, or else an abstract side of
the concept, which consequently has not yet been sufficiently thought
through. When the judgment is adequate, and the diremption complete, the
predicate, which posits the concept's concrete universality, must have the
same content as the subject, which presents its immediate singleness. The
whole concept has, according to the predicate, a reflected-through, logical
inwardness (\textit{in mente}): this is the concept's infinite ideality and
freedom; but, according to the subject, the whole concept is at the same
time posited in a corresponding, outer actuality (\textit{in re}), which
makes up its determinate reality. The diremption is thus driven to its
utmost boundary; but since subject and predicate have now also expressly won
the same content, the copula itself is developed into a full, real being, in
which the concept's extremes join together into concrete unity.

\parag{8}{The Logical Totality of the Judgment.}

Formal logic has further divided the judgments: a) with respect to the real
nexus of subject and predicate, into \emph{analytic} and \emph{synthetic};
b) with respect to the mutual agreement of the different sentences, into:
α) \emph{agreeing} (\textit{propositiones consentientes}) and β)
\emph{conflicting} (\textit{prop.\ oppositæ}) judgments, such that the
former again, according to their subdivisions, are either \emph{identical}
(\textit{pares, æquipollentes}), \emph{super=} and \emph{subordinate}
(\textit{subordinatæ}) or \emph{coordinate} (\textit{coordinatæ}), the
latter partly \emph{contradictory}, partly \emph{contrary}.

Since every judgment's content must always be explicated by the help of
another judgment, the speculative logic regards all \emph{single} judgments
as moments in \emph{the judgment's} --- i.e.\ the concept-sundering's ---
\emph{own logical totality}.

\anm{1} By examining how we come from a judgment's subject to its predicate,
it may appear partly that the subject according to its concept already
contains the predicate, partly that a special experience is presupposed in
order to join the predicate to its subject. In the first case the judgment
is called \emph{analytic}, because one needs only to make distinct to
oneself (analyze) the subject's content in order to find the predicate; in
the latter case, on the contrary, \emph{synthetic}, since predicate and
subject must each be sought for itself as separate concepts, which first
upon a logical deliberation let themselves be united. A closer consideration
shows, however, that it stands with this division as with all the rest; it
too has only a relative validity. The true analytic judgment is already as
such synthetic, since every predicate must join a new determination to its
subject; likewise every synthetic judgment is analytic, since the
predicate's union with the subject would lack all ground if it did not
consistently follow from the subject's peculiar character. Without the
analytic development, comprehension proper would be wholly wanting; without
the application of synthesis, thinking could make no advance (Kant:
Erleuterungs= und Erweiterungsurtheile). As the analytic judgments most
nearly express pure thinking, so the synthetic belong chiefly under
experience's field of vision; but the infinite unity of both lies at the
ground of true scientific cognition. As synthesis refers to representation
and the intellectual intuition, so analysis refers to reflection and the
pure insight; where synthesis most nearly presupposes difference and
opposition, analysis rests on identity and contradiction. The sentence: The
ground has a consequent, looks at first glance like a purely synthetic
judgment. The ground and the consequent are represented each for itself,
just as in the sentence: The door has a lock, one first represents to
oneself the door, then the lock, and finally grasps them both in their
union. But if the joining of subject and predicate is to express something
more than a mere fact, then the synthesis must proceed from an analysis.
That the latter sentence is hard to transform into a purely analytic
judgment lies precisely in its empirical character. The identity of subject
and predicate is a merely represented union, upon which the fundamental
proposition of contradiction finds only a \emph{critical application}. The
whole analysis aims solely at assuring the thinking subject that in the
concept of a door there is nothing that contradicts the concept of a lock.
But that no one who speaks of a door that has a conscience finds credence
proves, after all, that analytic thinking everywhere keeps a watchful eye on
experience's synthetic judgments. The first sentence: The ground has a
consequent, lets itself on the contrary be at once transformed into an
analytic judgment, just because it is of purely logical nature. There lies
no more in the ground than what comes into view in the consequent; for the
ground is ground only insofar as it has a consequent, and as the consequent
is, so must the ground also be. For pure insight the two concepts show
themselves as one concept; what the one is \textit{implicite}, the other
expresses \textit{explicite}. The consequent is thus ground for the ground,
and the ground a consequent of the consequent. But just as necessary is
nevertheless the essential difference of the two. The thorough contradiction
shows itself precisely in this, that analyzing thought each time comes to
posit their abstract identity when it wants to comprehend their pure
difference, and conversely must fix their abstract difference when it wants
to posit their pure identity. This contradiction reflection cannot solve
except by taking synthesis and intuition to its aid. It is first by means of
the intellectual intuition, the sound unprejudiced glance, that analysis's
abstract identity, integrated by synthesis's immediate difference, steps out
into determinate real opposition. It is therefore an error when many have
supposed that dialectical thinking should carry with it disdain for
intuitive cognition, since much rather, if it truly understands itself, it
must seek to penetrate and illuminate the real content of the represented
world.

\anm{2} As the sensing consciousness finds the multiplicity of things
immediately given, so also representational thinking comes to pass a host of
different judgments, and thereby to develop for itself a circle of sporadic
cognitions. It therefore becomes necessary to examine how these judgments
and cognitions stand to one another. For although consciousness in an
immediate way (--- faith, geniality ---) may possess an intellectual world
whose harmony and consistency can later be justified by the testing
reflection, still the uncultivated person, who receives an aggregate of
heterogeneous influences, intellectual as well as sensory, may often gather
in his consciousness such representations and concepts as, when brought
closer together, can by no means be reconciled in one circle of cognition.
From the subjective point of view, science must therefore be regarded as the
highest form in which a coherent view of life can be developed and
presented; and since logic is the general foundation of all scientific
culture, it cannot remain standing at a piecemeal apprehension of the
judgment's essence and significance, but must also demonstrate the different
judgments' inner nexus and totality.

As regards formal logic's division of \textit{propositiones consentientes},
the doctrine of the \emph{identical} judgments finds its determinate
application not only in rhetoric, where one and the same fundamental thought
is to be made intuitable from several different sides, but even in pure
dialectic, which in its category-analysis makes distinct the concept's
transitions and nuances. To determine more closely the fluid boundary that
separates the all-sided development of concept-identity from tautology
must, however, be left to the presenting subject's scientific tact.

% The Danish page's letterspaced run takes in the article: "de over= og
% underordnede" is spaced entire (compositor's lapse, logged in the
% transcription); the emphasis is carried here on the pair of adjectives.
With the \emph{super= and subordinate} judgments it stands as with the
super= and subordinate concepts, only that the criteria by which the logical
subordination is expressed in the judgment's form are determinately brought
out. It is then presupposed that the judgments in question must have either
predicate or subject in common. If the predicates are identical, then the
superordinate judgment has the superordinate subject; for the predicate's
extension is in this case measured by the subject (All men are mortal, all
animals are mortal); if the subjects are common, then the superordinate
judgment has the subordinate predicate; for the predicate's lower (more
concrete) concept must always be subsumed under the higher (All men have
sense-organs, all men have eyes).

As regards the \emph{opposed} judgments, it is presupposed that the talk
turns on one and the same object, wherefore formal logic also sets up the
rule that the \emph{contrary} judgments must have the same subject, the
\emph{contradictory} also the same predicate. These and similar
juxtapositions serve only to justify all the more the proposition that the
concept's totality can first come into view in a totality of connected
judgments.

\anm{3} What has already been said of the absolute concept --- namely, that
it takes up all relative concepts as moments and reflexes in its inner
totality --- can now also be applied to the judgment, regarded from the
standpoint of the higher logic. As it is the same thinking I and the same
judging understanding that utter themselves in all the individual's thoughts
and utterances, so it is also a proof of cognition's thoroughness, depth and
coherence that the one judgment lets itself be deduced from the other. While
formal logic therefore concerns itself most nearly with the judgments in
their sporadic plurality, the speculative treats of the judgment, according
to its essential unity, as that concept-sundering which with the concept's
development runs through its different, lower and higher stages. According
to the foregoing explication of the judgment's logical content, which began
with quality's (existence's) positive judgment, and ended with modality's
(the concept's) apodictic judgment, it has already been established that the
latter has taken the former up into itself, or, what is the same, that the
apodictic judgment is the positive judgment's truth; likewise every outward
separation of the formal and the real nexus between subject and predicate
must lose its significance when we hold fast the result won above, that the
judgment in all its forms is both analytic and synthetic. To assure oneself
that the categorical judgment presupposes analysis just as much as
synthesis, one needs only to make distinct to oneself the transition from
the general to the universal. That all men are mortal can first be known
with complete certainty when it is seen that the concept of mortality is
contained in the concept of man's natural existence. The separation of
concept and actuality, which in objective respect makes the judgment
hypothetical, makes it in subjective sense also problematic, and the
disjunction whereby a given genus-concept is sundered into its species is in
turn a presupposition from which the apodictic judgment results. For since
the apodictic judgment is to express the identity of actuality and concept,
it necessarily requires a disjunction. For no actual existence is
purposive, beautiful or good in general; but every thing that answers to
\emph{its} determination is always good in \emph{its} kind and in \emph{its}
species. Many unfair and unjust judgments, passed upon moral actions as well
as upon artistic and scientific performances, arise not merely from the fact
that the critics often find it most convenient to keep to an immediate
estimate, and therefore prefer to employ the assertoric mode of expression;
but the fundamental fault often lies in this, that the proper point of view
is displaced, because the category under which the object in question is to
be subsumed is missed altogether.

The categorical judgment receives its closer delimitation by means of the
disjunctive, and when this is faulty, the apodictic must also become false
and one-sided. But just for this reason it becomes necessary to distinguish
between the concept's one-sided and its all-sided expression. Although the
whole concept is present in each of its principal sides, it is still not
completed in any single moment. Hence it is also impossible to present the
whole concept in any isolated judgment whatsoever; for if the judgment has a
rich and concrete content, its moments must be \emph{analyzed} in a series
of different judgments; and every judgment whose content is limited and
one-sided must be synthetically integrated by another judgment, this again
by another, and so forth. --- The wish to express the highest concept in a
single judgment hangs together with the endeavor to preserve philosophy's
kernel in certain determinate principal propositions. True enough, a
contentful truth can be expressed in an adequate judgment; but even the most
pregnant proposition will at once sink down into an empty phrase when it is
torn out of the connection in which it fits, when it is no longer recognized
as the result of a preceding or the presupposition of a following
development. Hence the profoundest philosophy has always been the greatest
offense to so-called sound common sense; for it lies in the nature of the
case that the deeper the thoughts a philosophy contains, and the less
trouble one takes to press into its spirit, the more easily will its single
theses and utterances appear offensive, often even ridiculous.

\anm{4} The difference between the formal and the speculative logic's
apprehension of the judgment's relation to the concept can be still further
illuminated by looking to \emph{definition} and \emph{division}.

Insofar as the judgment's predicate expresses what the subject according to
its concept properly is, the adequate judgment can also be regarded as the
concept's definition. According to their immediate difference the concepts
can be separated and characterized just as exactly as the sensory objects.
It has therefore often been regarded as a proof of clear and acute thinking
to advance from definition to definition in the course of scientific
development. By the definition the concept is held fast as a magnitude
bounded in itself, firm and constant, which under the different
juxtapositions and turns is to keep the same value, and since every concept
must contain something peculiar whereby it is known from all others,
adequate definitions will always have their significance as contributions
toward clearing up meaning and securing linguistic usage. Nevertheless it is
chiefly the so-called exact sciences in which definition plays a
considerable role, just because thinking is here more mechanical than
organic. For philosophical cognition, on the contrary, such definitions have
a very subordinate worth. The ideal content that lies in the philosophical
root-words, such as: essence, ground, substance, and so on, can no more be
held fast in isolated definitions than the soul in an anatomized body. The
more thinking is lost in atomistic concepts and sharply bounded definitions,
the sooner will it approach the comfortless result that truth itself is
unnameable, incomprehensible and ungraspable; for no living spirit lets
itself be caught in a net of finite definitions. Formal logic has already
recognized that the \emph{genetic} definition stands above the
\emph{nominal}; but first in the stream of dialectic can the relative
concept's rise, culmination and decline be clearly mirrored, and the genetic
definition carried through. Since now every judgment contains a relative
definition, and since the one concept-moment must always be defined by the
other, it is evident from this that the absolute concept's content can be
presented only in a \emph{system of judgments}.

To distinguish itself from mere description, the nominal definition must
characterize the object by determining its concept, and delimit this by a
statement of \textit{genus proximum} and \textit{differentia specifica}.
But here a double regard is to be observed: \emph{partly} the concept must
be sought out in its immediate, empirical existence, and its different forms
and marks ``\emph{exposed},'' in order to find out its characteristic ---
and here it will matter whether one seizes an essential determination or
merely brings out an outer mark (synoptic and heuristic classification);
\emph{partly} singleness with its concrete peculiarity falls away, since the
definition does not present the concept's whole content, but only states its
general outline. The inner determinateness and separateness that the
definition renounces makes itself known in the logical division.

The transition from definition to division can take place by a double way:
a) partly the \textit{differentia specifica} stated in the definition can be
set against another specific difference falling under the same
\textit{genus}; for when man, for example, is said to be a \emph{rational}
animal, herein already lies that there are also animals that have no reason;
which again is the same as that the animal is \emph{either} rational
\emph{or} irrational, or: that all animals can be divided into the rational
and the irrational; b) partly the defined concept can be apprehended as a
new \textit{genus}, which within its own boundaries again yields new
\textit{differentiæ specificæ}, and thereby sunders itself into new species:
the rational animal is either religious or irreligious, i.e.\ men can be
divided into the religious and the irreligious. In the first case the
disjunctive judgment can immediately be set out as the division's logical
expression, so that the concept that is divided (\textit{divisum}) is
posited as subject, while the members of division (\textit{membra
dividentia}) find their place in the binomial or polynomial predicate. In
the latter case a subdivision (\textit{subdivisio}) arises, as one of the
predicate's members is again made subject in a new disjunctive judgment,
which then so far can be regarded as a consequence of the preceding. So long
as the disjunctive judgment is not apprehended as the division's logical
form, it still lacks the decisive determinateness whereby it distinguishes
itself from a merely problematic mode of expression. The sentence:
\textit{A} is either \textit{B} or \textit{C}, can namely be resolved into
two problematic judgments: \textit{A} is perhaps, when more closely
determined, \textit{B}, not at all \textit{C}; and: \textit{A} is perhaps
\textit{C}, not \textit{B}. Under this presupposition ``\emph{either}'' ---
``\emph{or}'' expresses an abstract exclusion, which rests in the
\textit{pricipium exclusi medii}.
% "pricipium" (for principium) is the Danish page's own misprint, logged in
% the transcription; carried as printed since it is quoted Latin matter.
To characterize the disjunctive judgment more closely, the sentence:
\textit{A} is either \textit{B} or \textit{C} or \textit{D}\ldots, must be
equivalent to the sentence: \textit{A} is both \textit{B} and \textit{C} and
\textit{D}\ldots The disjunction's: \emph{either} --- \emph{or} contains at
the same time the conjunction's ``\emph{both}'' --- ``\emph{and},'' from
which the division's relation to the concept is already evident. For if we
grasp the disjunction from the connecting as well as the excluding side,
herein lies that the same concept at once separates itself from itself into
a series of separate concepts, which mutually exclude one another, and
nevertheless comprises and binds them all within itself by repeating itself
in each single one. Thus it becomes necessary to seek out a medium in which
the exclusion and the connection can be thought united. \textit{A}, which in
one respect is \emph{both} \textit{B} \emph{and} \textit{C}, is in another
respect \emph{either} \textit{B} \emph{or} \textit{C}. Since both regards
lie in the disjunction's own essence, formal thinking must bring out a
universal determination in which the separateness comes into view, and thus
fixes the determinate \emph{ground of division} (\textit{fundamentum
divisionis}). It is, for example, a determination common to all triangles to
have 3 angles and 3 sides. But no actual triangle has angles and sides in
abstract generality, but always in determinate proportion. If now the
division is undertaken, it appears quite clearly that the concept's identity
at the same time grounds its difference, or in other words, that
\textit{diversi respectus} pass over into each other. For one cannot divide
the triangles into three-sided and three-sided; for here is an identity
without difference; nor into right-angled and isosceles; for here is a
difference that lacks its identity, or, as it is said, the ground of
division is not held fast. But when we say: the triangles are either
equilateral or non-equilateral, then it lies in the division's logical
nature that they are the one because they are the other. The contradiction
that hides itself in the disjunctive judgment --- that \textit{A} must be
either \textit{B} or \textit{C}, and yet is also posited as \textit{B} when
it is posited as \textit{A} --- finds its resolution in the division, which
sunders and binds the opposed concepts.

What \emph{the division} is from the empirical standpoint, the
concept-diremption or \emph{the primal division} is from the logical. All
the rules formal logic sets up for division point to this difference. That
the division shall determine the concept's extension by a statement of all
the genus's species means, in other words, that it shall exhaust the
genus-concept's content. But since every empirically characteristic mark can
always occasion new divisions or subdivisions, without insight into the
nature of the genus and the species being thereby increased, such a complete
enumeration is often as useless as it is difficult. ``Whether to the 67
species of parrots yet a dozen more can be found is indifferent for the
exhaustion of the genus-concept'' (Hegel). It is quite another matter that
detailed classifications may be necessary in order to complete empirical
knowledge; but they can have no properly logical worth. A division can often
be complete and correct without therefore being thorough, just as conversely
it can have a deeper significance without being logically correct (fantastic
and comic divisions). It is further conceded in formal logic that the
divided concept must be present whole and undivided in each of the members
of division; whereby the logical division distinguishes itself from the real
partition (\textit{partitio}). But from this it can also be seen that the
true division is grounded in the concept's primal division. When a sensory
object is divided, each part is only a piece. Hence it is also only the
elementary and lifeless things that can be taken apart and put together
without suffering harm, and that precisely because their concept is so
wholly external; the living organism, which expresses the concept's essence,
suffers on the contrary by every dismemberment, since it is contrary to the
concept's intellectual singleness to let itself be carved up. --- As
regards, now, the relation between the judgment and the division in general,
they both have the common determination of presenting the concept's
diremption. The criterion for such a diremption is always this, that the
concept is separated into two extremes, each of which for itself in its form
contains the whole concept; but while the judgment most nearly renders the
concept's single sides and reflexes, and so far lets itself be expressed in
a determinate sentence, the division on the contrary refers to different
concept-spheres, so that its content can be explicated only in a totality of
judgments. As each of the judgment's extremes, by the copula ``is,''
intimates a medium in which they meet, so it also lies in the division's
ground that the opposed concept-spheres must join together in a higher and
more concrete totality.

\secrule

\capitel{Third Chapter.}{The Validity of the Subjective Concept: Inference and Proof.}
\parag{9}{The Logical Requisites of Inference.}

The mutual relation of the judgments is expressed in the inference, that act
of thinking whereby a judgment's validity is derived from one or several
other judgments. According to formal logic there must be distinguished in
every inference a) one or more presuppositions, premises, b) a
conclusion-sentence (\textit{conclusum} or \textit{conclusio sensu minus
proprio}) and c) the whole act of thinking by which the conclusion-sentence
is derived from the premises (\textit{conclusio sensu proprio}). The
inference's simplest and yet at the same time complete form is the
categorical syllogism. The conclusion's subject, the minor concept
(\textit{terminus minor}), and its predicate, the major concept
(\textit{terminus major}), must in the syllogism each for itself be joined
with a third concept, which makes up the middle of the extremes
(\textit{terminus medius}). Of the two premises formed by this joining, that
in which the major concept is referred to the middle term is called
\textit{propositio major}; that, on the contrary, which unites the middle
term with the minor concept, \textit{propositio minor}. When the middle
term's relation and significance are distorted, paralogisms and sophisms
arise.

The speculative logic presents the inference as the extreme totalities'
dialectical transition into their medium, i.e.\ as a joining-together in
which the extremes are at once posited and annulled.

\anm{1} As expression of subjective thinking, the judgment proceeds from a
deliberation; every deliberation takes place in virtue of a judging; thus
the one judgment has its grounding in the other, or, what is equivalent, the
judgment has its truth in the inference. Every judgment that does not result
from an inference still stands as a postulate, which can be both true and
false, and for that reason both affirmed and denied. If the judgment is
regarded as answer to a question, then the inference contains the answering
of a double question, and thereby expresses the judgment's ground. If it is
thus asked: ``What is God?'' then the preliminary answer: ``God is a
spirit,'' is to be regarded as an adequate judgment. But if the judgment is
to be grounded, then the question resolves itself into two others: ``What is
God, since he is called spirit, and what is spirit, since it is predicated
of God?'' There is thus sought a third concept, which expresses the identity
of the subject's and the predicate's unity and difference --- an advance
whereby the moment of separateness first comes into its full right.

When the difference between subject and predicate does not come forward with
such independence that the concept grounding the unity of both needs to show
its relation to each of the members in a separate judgment, the conclusion
can be derived from one premise, and the inference is characterized as
immediate. Formal logic distinguishes several kinds of immediate inferences,
and sets up rules for their application. We thus obtain: a) the inference of
identity (\textit{ratiocinium ad æquipollens}), b) the inference of
subordination (\textit{ratiocinium ab universali ad particulare}), c) the
inference of opposition (\textit{ratiocinium ad oppositum}) and finally d)
the inference of conversion (\textit{ratiocinium per conversionem}). But
since none of all these inferences leads to any deeper insight into the
essence of the concept or of the judgment, they can only be used as simple
exercises of the understanding.

\anm{2} It is self-evident that the so-called categorical syllogism contains
all the requisites of inference, since it is first in it that the
conclusion's subject and predicate are comprehended in their independence,
and yet at the same time find the separate concept that makes up the union
of both. Everything thus rests essentially on the middle concept's
significance and relation to the extremes. Formal logic has for that reason
set up the following rules:

\begin{itemize}
\item[1)] \textit{Ex meris negativis nihil sequitur}; quite naturally,
  since the negation, when it holds for both premises, must cut off every
  connection between the extremes and the middle concept.
\item[2)] \textit{Ex meris particularibus nihil sequitur}; for two
  particular premises make the connection of the extremes and the middle
  concept doubtful.
\item[3)] \textit{Conclusio sequi debet partem debiliorem præmissarum}. For
  nothing can be posited in the conclusion that is not \textit{implicite}
  presupposed in the juxtaposition of the two premises.
\end{itemize}

\anm{3} The ambiguity that \textit{diversi respectus} contain, in that they
both produce and annul the contradiction, shows itself most clearly of all
in syllogistic. For since every thing can be regarded from more than one
side, one can always find for a conclusion's subject and predicate more than
one middle concept, just as conversely the middle concept and its relations
can in turn be apprehended from several different points of view
(Leibnitz's combinatorial analysis --- Plouquet's calculus). As subjective
reflection hereby develops itself into formal freedom, and uses this freedom
arbitrarily to set truth's semblance in truth's place, syllogistic passes
over into sophistry. The sophisms of the ancients do not, however, all bear
the same stamp; for while some, as for example the so-called
\textit{Elenchus in dictione}:

\begin{center}
\textit{mus rodit caseum,}\\
\textit{mus est syllaba,}\\
\textit{ergo: syllaba rodit caseum.}
\end{center}

\noindent are so crass that they must be regarded as logical
``Kjældermænd'' (flat puns), others on the contrary (--- ψευδόμενος,
σωρείτης, φαλαρκός ---) conceal a dialectical sting, and thereby form a kind
% φαλαρκός, with ρκ where the Greek wants κρ, is the Danish page's own wrong
% setting, logged in the transcription; copied verbatim.
of holding-points for acute thinking. It is the contradictions of the
represented world and of the so-called sound understanding that have given
the first impulses to philosophical movements of thought; but so long as a
more determinate method was still lacking, the deeper reflection had to keep
to sporadic παράδοξα, and make the problems intuitable by giving them a
piquant mask. Only as vehicles for ingenious plays of thought do the
sophisms have a relative significance; when on the contrary they are
cultivated for their own sake, they testify to philosophy's vanity and
self-dissolution. Even the theory of the refutation of the sophisms is
without logical interest; for he who has first learned to think consistently
and thoroughly will in every case know how to unmask the false, and a
reasoning that departs altogether from all sound thinking does not even
deserve to be refuted. Least of all can the rules for such refutations be
treated in the doctrine of inference; for the syllogism is only the general
form under which the sophists insinuate the most heterogeneous distortions.
Aristotle's treatise περὶ σοφιστιχῶν ἐλέγχων therefore yields only a meagre
profit for logic, but is nonetheless a document remarkable for the history
of philosophy, insofar as it contains a testimony to the enormous
thought-empiricism, the wealth of reflection-experiments, in which Greek
philosophy had to try itself before the systematic logician could appear.

\anm{4} As the judgment is the subjective thought-form that corresponds to
the objective concept-sundering, so also the syllogism is to be regarded as
a silhouette, drawn by the understanding, of that infinite circuit in which
the concept's separated extremes seek their concrete unity. The difference
between the formal and the speculative apprehension of the inference's
essence and significance rests then chiefly on this, that the logic of the
understanding sets up the premises as given sentences, whose middle term
must be taken as it is found, while the logic of reason on the contrary
reduces the premises to opposed concept-totalities, whose medium develops
itself out of their dialectical relation; in the logic of the understanding
the conclusion (i.e.\ the \textit{conclusum}) is a consequent sentence, a
result the thinking subject itself draws out; according to the logic of
reason it is on the contrary the sundered totalities' mutual transition, in
which they integrate each other --- the joining-together in which they both
find their subsistence. In the formal inference of the understanding, the
syllogism, the single members are set out in an outward succession; in the
true inference of reason they are simultaneous. Thus the coherence of the
universe consists in an inference: the sun draws the planets to its centre;
every planet in turn flees the attracting sun, and the medium of the
centrifugal and the centripetal force is the elliptical movement. The animal
development of life goes on in the form of inference; for the organism is
the medium between the life-moments' undivided unity in the soul's inner and
the functions' multiplicity in the bodily outer. Freedom is an inference: it
presupposes itself in its endowment, actualizes itself in its striving,
possesses itself in its result, and the means is always medium between
purpose and execution. The examples here adduced certainly do not further
explain either the universe, life or freedom, but they serve nonetheless to
make intuitable the truth that the inference is the objective form of
reason, and to make intelligible the Hegelian utterance: ``der Schluß ist
nicht nur vernünftig, sondern alles Vernünftige ist ein Schluß.''

\parag{10}{The Formal Validity of Inference.}

In order to express its true logical content, the inference must run through
several stages of development, and the forms of inference can therefore be
divided: a) according to \emph{quality}, into the inference of separateness,
of singleness and of universality, such that the \textit{terminus medius},
which here expresses only an abstract concept-determination, runs through a
threefold schema: \textit{E---S---A}, \textit{S---E---A} and
\textit{E---A---S}, in order to mediate the extremes; b) according to
\emph{quantity}, into the inference of allness, induction and analogy;
reflection, which empirically gathers and separates the concept's essential
universality and existential singleness, strives at this stage out beyond
the rule toward the immanent law; c) according to \emph{relation}, into the
categorical, hypothetical and disjunctive inference, whereby the concept's
ideal sundering and actuality's real sundering reflect each other with such
necessity that the totality comes into view in all forms.

According to the speculative logic, the inference's absolute truth is to be
explicated in a system of three inferences.

\anm{1} The so-called figures of inference, of which Aristotle set up the
first three, and later logicians (Galen) added a fourth, are distinguished
in formal logic according to the place the middle term occupies in the two
premises. Thus the first figure is characterized by the middle term's being
subject in the major premise and predicate in the minor; according to the
second figure the middle term is predicate, according to the third figure
subject, in both premises. The fourth figure, in which the middle term is
subject in the minor premise and predicate in the major, forms a superfluous
contrast to the first. To each of these figures determinate rules are then
attached: as regards the first figure, the major premise must be universal,
the minor affirmative, and the whole subsumption follow the general norm:
\textit{nota notæ est nota rei ipsius}. In the second figure, whose major
premise must likewise be universal, the one premise is affirmative and the
other negative; the third figure demands that one of the premises be
universal, and the minor premise affirmative; the fourth figure is
artificially produced by various conversions. This schematizing shows that
each of the concept's principal sides must come forward as \textit{terminus
medius} in order that the joining of the extremes may be fully grounded;
but, isolated from their logical content and mutual connection, these
abstract forms of inference have no deeper significance. By reducing the one
figure to the other, formal logic awakens a presentiment of their inner
connection.

% The Danish prints this remark as "Anm. 1." a SECOND time — § 10's remarks
% run 1, 1, 3, 4, 5 in the original (see the transcription's notes at pp. 75
% and 79). The printed numbering is mirrored here.
\anm{1} The inference of quality has the mark that it always keeps to the
immediate difference and unity of the concept-determinations. What is only a
single cannot as such be a universal. In order that two such abstract outer
members may be united, a middle determination is required, which in one
respect coincides with singleness, in another respect with universality;
this middle determination is the \emph{separate}. ``Bravery is a virtue,
Alexander has shown bravery; Alexander is virtuous.'' The qualitative
inference thus begins with the first figure: \textit{E---S---A}. That this
inference is very imperfect appears upon closer consideration. The
conclusion, whose validity was to be derived from the premises, itself
conditions their truth. For if the conclusion is merely to declare that
everyone who is brave is also, without any regard to his other qualities,
virtuous, then the conclusion is only a tautological repetition of the
\textit{propositio major}, and thinking makes no advance whatever. If on the
contrary it is to be expressly brought out that a single, determinate
individual is virtuous on account of \emph{his} bravery, then the inference
is false, since the conclusion in this case contains more than can be
derived from the premises. By bringing out another quality in the individual
concerned, one can therefore also come to an entirely opposite result.
``The sensuous nature is not virtuous, Alexander has a sensuous nature,
therefore he is not virtuous.'' Whether a separate and a universal
concept-determination let themselves be united or not depends then upon the
concrete single existence, which according to its nature always has several
different qualities. Singleness must accordingly be laid at the ground of
the joining of the separate and the universal, whereby the transition takes
place to the second schema (i.e.\ formal logic's third figure):
\textit{S---E---A}: ``Man is rational, man is mortal; a mortal is
rational.'' The separate determination ``mortal'' and the universal
determination ``rational'' are found united in every human individual. More
closely regarded, this conclusion likewise expresses no proper inference.
The singleness or the individual must simply be sought out, and when it is
found, then the separate and the universal determinations that inhere in it
are at once discovered by analyzing thought. Hence it is also in itself
quite indifferent which of the two extremes one will in this case make
\textit{terminus major}; the conclusion: ``a mortal is rational,'' can
without further ado be turned about: ``a rational being is mortal.'' The
proper advance consists consequently in the middle concept's contingency
being expressly brought out. --- In order to form the true middle between
the separate and the universal, the single must itself be determined as
separate; this can happen only in virtue of universality. With this
presupposition the transition is made to the third schema (formal logic's
second figure): \textit{E---A---S}. Insofar as this form of inference is
regarded in its connection with the two preceding, both its premises can be
said to be formally grounded. For the \textit{propositio major}:
\textit{S---A} stands as conclusion in the second schema, and the
\textit{propositio minor}: \textit{E---A} results from the first.
Nevertheless the middle concept is here too insufficient, since the
qualitative universality falls outside its singleness and separateness. This
defect comes into view in the inference's ending in a negative result: ``No
man is four-legged, the lion is four-legged; the lion is not man.''

It stands with the qualitative inference as with the qualitative judgment;
its moments can, according to their unity and difference, be reduced to
tautological sentences and superficial contradictions. \emph{The single is
only single}, insofar as it is held fast one-sidedly and immediately.
``This apple'' is, for example, a single, for one is to think neither of its
manifold qualities, nor of its likeness and unlikeness with others. But
every other apple is again, taken by itself, a single, a ``this,'' which is
immediately pointed out. The consideration of the existing things thus leads
us from the single to the single, without any separateness or universality.
This tautology shows itself at the same time as a superficial contradiction:
by repeating itself, \emph{the single becomes universal}, and since this
repetition cannot be thought without the difference between the several
singlenesses coming into view, \emph{each single must immediately be
determined as a separate}. By analyzing the concrete singleness one can
separate it into a plurality of universal determinations. The single apple
has thus a certain fragrance, taste, form, weight and so on, and when all
the universal determinations are separated, the thing's singleness vanishes
at the same time. As the nominalist annuls \textit{universalia} in
\textit{res}, so the realist dissolves \textit{res} into its
\textit{universalia}. If now each \textit{universale} is held fast for
itself, we have again both a tautology and a contradiction. ``Fragrance is
fragrance, without regard to that which is fragrant; taste is taste, without
regard to that which tastes:'' \emph{the universal is universal}. But as the
fragrance, the taste, the color and the weight are each in particular
posited as qualitative universalities or universal qualities, each
\textit{universale} becomes itself single: \emph{the universal is a
single}. Although now every universal determination must be observed for
itself, it still cannot be defined except by the help of another universal
determination, so that the difference between \textit{universalia} is by no
means indifferent: \emph{the universal is a separate}. It is further evident
from this that the all-sided determination of separateness is the
inference's aim --- an aim the qualitative inference cannot realize. In
order to ground the true unity of the single and the universal, separateness
must at once be identical with both extremes, and yet at the same time
different from them both. But now it lies in quality's own concept that the
separate, when it is held against the single, becomes itself a single, and
when it is set against the universal, is transformed into an abstract
universality. Hence the \textit{terminus medius} alternates between the
three different extremities: \textit{S}, \textit{E} and \textit{A}. Thus one
can, for example, lay the genus, the separate, at the ground of a joining of
the individuals and the universe. In its difference from individual and
universe the genus is then hypostatized and personified. But just as
immediately one can then also turn the whole form of inference about, and
declare that the genus exists only in the single individuals, that the
universal must not be thought separate from the single, and thus make the
individual \textit{terminus medius} for genus and universe. If it must now
further be conceded that the individuals unceasingly dissolve and pass away,
then one needs only to invert the form of inference once more, and lay a
universal quality (water, air, fire) at the ground of a conclusion of
individuals and genus. Since it is now an impossibility to state how a
universal element of nature could have sundered itself into separate
existences, this inference gives no positive result whatever, but expresses
only the negative truth that the separate genera and their individuals could
not exist if they did not have a common natural medium. By leaping from the
one middle concept to the other, naturalism always confuses the condition
with the creating ground, and the philosophy of atheism characterizes itself
by thinking's never going out beyond the qualitative inference.

\anm{3} The formalism in which the inference of quality ends finds its
purest expression in the inference of quantity. According to its immediate
quantity-content, the inference can be reduced to the mathematical axiom:
two magnitudes, each of which is equal to one and the same third, are equal
to each other; \textit{a} = \textit{b}, \textit{b} = \textit{c}, therefore
\textit{a} = \textit{c}. But since all the \textit{termini} here have the
selfsame significance, there can in the logical sense be talk neither of
concept, judgment nor inference. The inference of quantity must for that
reason be further developed and determined as the \emph{comprehensive}
inference, in which the universal and the single hold themselves with
relative independence over against each other, and therefore seek a more
concrete \textit{terminus medius}. For the rest, the inference of quantity
runs through the same schemata as that of quality. In allness the universal
and the single come immediately into view with their unity and difference.
Since allness is the collective universality, singleness becomes by its own
repetition itself universalized; but since every single in turn has its
independent significance in the allness, universality is at the same time
based upon singlenesses; the medium or the separate is in this case nothing
other than the universal's and the single's existential determinateness.
Hence the comprehensive inference begins, as \emph{inference of allness},
with the first schema: \textit{E---S---A}. ``\emph{All} (single) men are
mortal, Caius is \emph{a single} man; therefore C.\ is mortal.'' The
conclusion is in this form too a tautological repetition of the major
premise, and the inference of allness is for that reason only to be regarded
as a transition, so that everything depends upon an actual union of the
separate and the universal, whose ground is to be sought in singleness. The
comprehensive inference hereby enters into the second schema:
\textit{S---E---A}, and is determined as \emph{induction}. According to
formal logic, inferences of induction are characterized thus: ``one infers
that a predicate that belongs to \emph{a determinate} number of subjects of
a certain \textit{genus} must hold for \emph{all} subjects belonging under
this \textit{genus};'' or: ``What holds of certain species and individuals
that are subsumed under a universal genus-concept must also hold of the
whole genus.'' But from this point of view the inference loses its essential
significance; for if all possible cases are not enumerated, it is
incomplete, and when they are all adduced, one obtains an addition instead
of an inference. If induction's logical worth is to be acknowledged, then
the accent must necessarily fall upon the \textit{terminus medius}, which
here presents the developed allness. More closely inspected, allness now
offers a double side; for partly it shows itself to thought as the essential
universality, partly it is apprehended by immediate observation as the
common determination of the singlenesses. The essence, the universal, exists
only in the individuals; one cannot therefore cognize the genus without
having first observed all its species and individuals; but every following
individual is always in certain ways different from the preceding, and the
number of the singlenesses so great that one can never finish the
observation. Cognition's a priori moment is still hampered by the
preponderance of the a posteriori, and the result is a combination of
thinking and sensation, which lies at the ground of all experience. By
posing a task that cannot be solved, induction marks the logical transition
to the inference of analogy. Instead of seeking the universal in the sum of
the singlenesses and losing itself in an endless approximation, thinking now
posits a given exemplar as immediate expression of the genus-concept and its
specific determinations, whereby the inference's schema is changed to:
\textit{E---A---S}. If two subjects have a certain number of predicates in
common, one infers according to the formal analogy that they must have yet
more. The earth and the moon have, as planets, several common properties;
since now the earth is also inhabited, it seems to follow from this that the
moon must be so too. The moon (\textit{E}) is thus ascribed, on account of
its likeness with the earth (\textit{A}, insofar namely as the earth in this
case represents the planet's concept), the particular determination
(\textit{S}) of being inhabited. The general concept, which fled through
induction's endless progress, is now immediately held fast as the inference
of analogy's determinate foundation; but since its appearance is quite
empirical, it still remains doubtful how many of the particular
determinations belong to the genus, and consequently can be found again in
all its species and individuals, and how many on the contrary belong to the
single exemplar. Even if natural science had not pointed out phenomena that
conflict with the hypothesis of moon-dwellers, it would still always be
problematic whether it belonged to the earth as a single individual globe,
or as the planet-concept's existential expression, to be inhabited by living
beings. As astrology shows how one can misuse induction and analogy to
introduce fantastic dream-images into actuality's realm, so Christophobia
proves on the contrary, insofar as it gives itself a philosophical tinge,
that the same forms of inference can be used when the aim is to transform
the historical truth into a dream-image. It is this hovering between the
immediate experience and the free thinking, between the outer actuality and
the essential concept, that characterizes the comprehensive inferences, and
occasions their great ambiguity. Reflection can fall into the
leading-strings of sensation and empiricism; but the immediate glance of
thought is again able to penetrate the empirical surface and gather the
sporadic moments in their essential ground. The rule that develops itself in
order to become law, the circuit and the endless progress, repetition and
alternation, belong under this sphere; hope itself is an inference of
induction. Hence the inference of induction and analogy extends its logical
dominion all the way from the prosaic probability and the empty surmise,
through the deep presentiment, to the genial divination.

\anm{4} The opposition between the thought and the existing concept,
together with the concrete unity of both, is developed in the inference of
relation, whose mark is that each of the three \textit{termini} itself
expresses the whole concept-totality. The difference between the extremes
therefore seems at first glance to concern not so much the content as the
form; for every single concept, individual, is the general concept's own
actuality, and the concrete general concept is in turn the actual
individual's totalized essentiality. But, more closely regarded, the formal
difference is also to be recognized as a difference in the content itself.
For as immediate single existence the individual can have many indifferent
and accidental qualities that fall outside its concept; it can moreover find
itself in such a state that the essential and original determinations are in
many ways hidden and hampered, so that the actual exemplar by no means
expresses the whole concept-content. A similar one-sidedness obtains with
respect to the general concept, whose separate determinations can appear
only in and by means of the individual. If the sensuous did not exist, then
sensuousness could not be either; if the rational did not exist, the thought
of a universal rationality would be a meaningless abstraction. The true
unity of the single and the universal is consequently to be determined as
both extremes' infinite relation, i.e.\ such a joining of both outer members
that the individual itself is developed to its concept, whereby
\textit{universalia in re} are transformed into \textit{universalia in
mente}; while the concept is at the same time individualized, and
\textit{universalia in mente} converted into \textit{universalia in re}. As
a natural individual without upbringing, man cannot reach his determination,
because he lacks the stamp of the universally human, and thus, logically
speaking, does not answer to his concept; but just as little can a blind,
habitual submission to the laws of duty and morality satisfy human nature,
which will always demand an independent appropriation of what is handed
down, as the free single will's imperishable right. True upbringing is an
inference of relation, which has the natural individuality and the moral
laws for its extremes, but the living, personal unity of both for its
\textit{terminus medius} (--- false originality, formal school, true
originality; --- mystical self-annihilation, mystical self-deification;
religious personality ---).

Since the concept's essence in relation to the existing individuals must
always be the overreaching, insofar as it expresses the things' category,
the category or the separate concept forms the medium between singleness and
the abstract universality. Hence relation begins with the \emph{categorical}
inference after the schema: \textit{E---S---A}. In syllogistic form the
categorical inference cannot well be distinguished from the inference of
allness, since the syllogism's middle term is always simply inserted without
developing itself out of the connection itself. What on the contrary
distinguishes the categorical inference above the general one is the insight
into the concept's essence, which makes every enumeration of the special
cases and empirical appearances superfluous. He who has once fully
recognized which fundamental determinations belong together in order to
constitute human nature can always be certain of finding each of these
again, either \textit{implicite eller explicite}, in the single individuals,
% "eller" is antiqua in the Danish and stands inside the italic run, per the
% transcription's note; carried as printed.
without needing an endlessly continued experience. If such categorical
certainty were not given, then every question of principle would be
meaningless, and science itself an impossibility. By means of the
categorical inference the concept's apriority is first fully acknowledged.
But as the categorical inference dissolves the existing things into their
categories, it presents at the same time a one-sided concept-realism
(\textit{universalia ante rem}), which must be corrected by a higher form of
inference. That no thing can contradict its category without changing name
and existence --- from this it does not yet follow that the thing's
individual determination should immediately coincide with the concept's
universality; because possibility's extension is determinate, its content is
not yet for that reason exhausted. In order therefore to satisfy
individuation and the higher realism (\textit{univ.\ in re}), singleness
must itself be laid at the ground. The a posteriori (experience and fact)
thereby steps up anew against the concept, and the mutual relation of both
is developed in the \emph{hypothetical} inference of the schema:
\textit{S---E---A}. Since the middle concept is reduced to singleness and
fact, and consequently must be given immediately, the hypothetical inference
can express its logical content in the syllogism's form. The major premise
then becomes a hypothetical judgment, and formal logic now prescribes a
double method of inference: from the antecedent's correctness one infers by
\textit{modus ponens} to the consequent's validity, after the rule:
\textit{posita conditione, ponitur etiam conditionatum}; from the
consequent's incorrectness one infers on the contrary by \textit{modus
tollens} to the antecedent's rejectability, according to the rule:
\textit{sublato conditionato, tollitur etiam conditio}.

Misled by a grammatical regard, formal logic has, however, committed the
error of confusing \textit{conditio} with \textit{conditionatum}. In a
judgment such as this: ``If \textit{F} has seen this play, it must have been
performed,'' the consequent evidently declares the proper condition, the
antecedent on the contrary the conditioned. The general relation between
actuality and concept is presented in the \textit{propositio major}, which
expresses the a priori law of thought without regard to any determinate
actuality. The minor premise thereby acquires, in consequence, a double
significance: partly the condition can be negated, in which case the
conditioned actuality is to be regarded as impossible (\textit{a non posse
esse ad non esse valet consequentia}); partly the single fact can show
itself so unambiguous that the presence of its conditions can be derived
from it (\textit{ab esse ad posse esse valet consequentia}). The application
of the hypothetical inference serves to make intuitable the true interaction
between thinking and experience. On this relation rests philosophy's formal
cultivation and real advance. A philosophical school that, for example,
overvalues the logical can, as regards principle, always be referred to the
% The quotation below opens with `` and never closes — the Danish p. 86 sets
% no closing mark (printer's defect, logged in the transcription); mirrored.
hypothetical major premise: ``If \textit{A} is a philosophical problem, it
must be solvable by pure dialectic. So long as such a school is in process
of forming itself, its adherents will, in the choice and treatment of the
philosophical tasks, keep to the presupposition's infallibility, and by
application of \textit{modus ponens} volatilize every problem whose solution
oversteps the boundaries of formal dialectic. ``If evil is a problem, it
must find a dialectical solution; now it is a philosophical problem,
therefore it has no other significance than that which formal dialectic can
allow it.'' But when the facts speak too loudly against such a reduction,
philosophy is at last compelled to take a step forward, turn the inference
about, and avow \textit{modus tollens}: ``If formal dialectic yields
philosophy's absolute method, then it must be able to solve the problem of
evil; now it cannot solve it, therefore the method is insufficient.'' If the
form of inference is apprehended from a general point of view, then the
one-sidedness lies in this, that the major premise states only an abstract
necessity, while the minor introduces a pure contingency; hence the
conclusion still lacks the true grounding. This defect the
\emph{disjunctive} inference remedies by positing its middle concept as the
real possibility of the actual existential forms, after the schema:
\textit{E---A---S}. This schema may indeed seem to deviate from the
disjunctive inference's usual form: \textit{A} is either \textit{E} or
\textit{E'}, or \textit{E''}; now \textit{A} is posited as \textit{E},
therefore it is neither \textit{E'} nor \textit{E''}; but according to the
significance ascribed above (\S\,8, Remark 4) to the disjunction, the major
premise, according to its concept, is to be presented thus: \textit{E},
\textit{E'}, \textit{E''} are, each for itself, to be subsumed under
\textit{A}; now \textit{A} is determined as \textit{E}, therefore it is in
this \emph{particular} (\textit{S}) case neither \textit{E'} nor
\textit{E''}. The general concept (\textit{A}) has now so concrete a content
that its real possibility preforms the outer actuality. The conditions for
its stepping out into particular existential forms are indeed still fetched
from without; but the existence that yields these conditions has itself
proceeded from the absolute concept's own ground. According to the
possibilities' inner sundering, \textit{A} is both \textit{E}, \textit{E'}
and \textit{E''}; but according to actuality's exclusive determinateness it
cannot be \textit{E'} and \textit{E''} when it is \textit{E}. If the outer
disjunction were not posited, then possibility would remain eternally
unactual; if the inner disjunction were only an empty abstraction, the
exclusive actuality would be impossible. As the disjunctive inference thus
yields the highest norm for the concept's formal development, it confirms at
the same time the truth that the absolute concept would have no validity if
it did not at once exist both \emph{inwardly}, in the all-intelligence's
infinite freedom, as a realm of sundered possibilities and individualized
potencies, and \emph{outwardly}, in the universe's immediate actuality,
whose outer totality stamps out the real disjunction. In essentiality
mediation holds; in actuality, \textit{principium exclusi medii}; the true
unity of theory and practice will therefore always rest upon a thorough
synthesis of mediation and exclusion.

\anm{5} The number of the inferences is, with respect to their separate
content, as infinite as the multiplicity of thoughts and things. It is
consequently an impossibility to enumerate all the cases in which they can
cross one another and pass over into one another. But in all different
sunderings and joinings the general ground-norm nevertheless holds, that
each of the three \textit{termini} is to be posited both as extreme and as
medium. Hence the higher logic teaches that every inference consists in a
system of three inferences. ``Many disagreements and confusions in science
have their ground in one's keeping to the single inference instead of to the
whole system. Hereby arise different, indeed opposite assertions, of which
each is correct insofar as it is a side of the truth, but incorrect insofar
as its truth is one-sided. When one says that man's spiritual nature
develops itself out of the sensuous, or that the body develops the soul, one
infers according to \textit{E---S---A}. If one says that the sensuous and
the spiritual are inseparably bound together in the individual, and that
neither of the two has priority over the other, then it is
\textit{S---E---A}. And if one finally assumes that it is the spiritual
nature that lies at the ground of the development of the sensuous, or that
it is the soul that forms its body, one infers according to
\textit{E---A---S}. But the whole truth is only in the complete system of
all three inferences. The first inference is materialism's, the second is
empiricism's, the third is spiritualism's. But the whole system is
idealism's.'' Heiberg: Den specul.\ Logik, Pag.\ 99.

\parag{11}{The Real Validity of Inference: Proof and Method.}

Since the formal syllogistic satisfies only the empty reflection, without
any regard to existence's real content, this defect must in turn be made
good by a series of immediate presuppositions being taken up outright as
given points of attachment for reasoning thought. Knowledge proper, or
scientific cognition, will thus rest upon the mediate or resulting unity of
thinking and being, which is to be sought in \emph{the proof}. But so long
as the proof is conditioned by so external a combination of concept and
actuality, the scientific method can be driven from the one extreme to the
other. Dogmatism, which builds upon unproven presuppositions, is supplanted
by criticism, which shakes all positive supports and overturns the thetic
edifice of doctrine. In order to escape the skeptical confusion, the higher
knowledge, which has not yet found its true method, must then develop
itself either into a formal idealism, which absorbs existence's content in
an immanent demonstration, or into a realism of feeling, which exchanges
the scientific form for a lyrical rhetoric, and relies upon immediate
evidence.

The true philosophical method can consist neither in a merely poetic-genial
unveiling of existence's absolute content, nor in a merely logical formal
development of the universe's concept-moments, but must justify itself as a
totality-dialectic, in which existence's principles show themselves with
their relative significance, and all life's potencies stir according to
their peculiar validity.

\anm{1} Aristotle (\textit{Analytic.\ post.\ c.\ 1---3}) places the
difference between syllogism and proof in this, that the syllogism in and
for itself gives no real cognition (συλλογιςμὸς οὐ ποιήσει ἐπιστήμην),
% The medially standing final sigma in συλλογιςμὸς (twice) is the Danish
% page's own wrong sort, logged in the transcription; copied verbatim.
while the proof (ἀπόδειξις) on the contrary always carries with it the
certainty that the matter must in itself stand as we think it, wherefore the
proof can also be called a cognition-syllogism (συλλογιςμὸς ἐπιστημονικός).
Nevertheless he seeks to refute the view that all cognition of truth should
be dependent upon syllogistic proofs. By positing the demonstration as
condition for all knowledge, one must either, as some have wished, deny
science's possibility, since every inference contains premises that must in
turn be proved by new inferences, so that the whole conduct of proof loses
itself in an endless regress; or else, what others have attempted, derive
the premises from their own conclusions, whereby the so-called
\textit{circulus in probando} arises, and the proof comes to hover in the
air. As there are, then, according to Aristotle, truths that can and must be
demonstrated, so there are also truths that require no proof, since they
themselves yield the original supporting-points for all demonstration:
ἀναγκη καὶ τὴν ἀποδεικτικὴν ἐπιστήμην ἐξ ἀληθῶν τ' εἶναι καὶ πρώτων, καὶ
ἀμέσων, καὶ γνωριμοτέρων καὶ προτέρων, καὶ αἰτίων τοῦ συμπεράςματος. οὑτος
γαρ ἔσονται καὶ αἱ ἀρχαὶ οἰκεῖαι τοῦ δεικνυμένου (\textit{Anal.\ post.\ c.\
2}). It is chiefly the assumption of these many ἀρχαί, or principles and
grounds of cognition, that separates Aristotle from Plato. For while Plato
aimed especially at explaining the multiplicity of things in virtue of the
universal concepts' (the ideas') apriority, Aristotle fixed his gaze also
upon the world of experience, in order to find in it the sporadic points of
departure, and to lead the existing single concepts (αἱ πρώται οὐσίαι) back
to their universal fundamental concept. Although now every science has its
separate presuppositions, still in all of them, even in philosophy, there
must both be something from which one proves, and also something about which
the conduct of proof turns; the Aristotelian doctrine of thought can
therefore be said to have established philosophy's dogmatic method.

\anm{2} After dogmatism in the history of philosophy had run through various
stages of development, it was in modern times driven to its peak by Christ.\
Wolf, who methodically cultivated the so-called \textit{methodus
mathematica}. Since every presupposition in this method coincides with a
definition that is to express the concept's essential determinations,
analysis is sufficient to bring out the result already anticipated in the
definition itself; whereby the whole demonstration acquires a semblance of
apodictic infallibility. ``The divine Being must, according to its concept,
possess all perfections; being is a perfection, therefore God must exist.''
--- The insufficiency of this dogmatic method was demonstrated by Kant
(jfr.\ Die Disciplin der reinen Vernunft im dogmatischen Gebrauche; Kritik
d.\ r.\ V.; Methodenlehre 1 Haupst.). While Kant concedes to mathematics
that it gives a brilliant example of pure reason's self-development without
experience's help, he at the same time denies the mathematical method's
applicability in philosophy. The methodical difference is grounded in the
peculiar nature of the two sciences. For while philosophical cognition
considers the separate in the universal, i.e.\ holds fast universality
\textit{in abstracto}, the mathematical on the contrary posits the universal
in the separate, and thus possesses it \textit{in concreto}; what ``der
Vernunftgebrauch durch Construction der Begriffe'' is in mathematics, ``der
Vernunftgebrauch nach Begriffen'' is in philosophy. Mathematics' definitions
are always infallible, since they are brought forth \textit{a priori}
together with their corresponding objects; while the philosophical
definitions on the contrary presuppose a content given in actuality, whose
determinations cannot be immediately delimited. Mathematics, which
constructs its concepts in a priori synthesis, can employ the demonstration;
whereas philosophy, whose universal concepts must always refer to
``possible experience,'' cannot attain to any speculative synthetic
judgment, and therefore cannot demonstrate its way forward either. But
although Kant thus sought to drive all ``dogmata'' from philosophy's
precincts, he still leaves behind, in his ``discursive Grundsätze'' and his
dialectical ``Antithetik,'' a remnant of dogmatism and demonstration.
% The Danish paragraph ends with a comma where a full stop belongs (printer's
% defect, logged in the transcription); normal punctuation is supplied here.

\anm{3} In his Wissenschaftslehre Fichte set himself the task of reconciling
dogmatism with criticism, and insofar as dogmatism refers to a ``Ding an
sich,'' it may be said to be thoroughly overcome in his transcendental
idealism; on the other hand he approved the demonstration, and developed it
to its highest stringency. As the first unconditioned principle: \textit{A}
= \textit{A} is annulled in the thinking I's identity with itself, so the
second: \textit{A} is not = --- \textit{A} expresses the I's difference from
the not-I. \emph{Identity} and \emph{opposition} thus give themselves as
expressions of I-hood's original fundamental acts. The third principal
proposition: \textit{A} and --- \textit{A}, I and not-I, mutually
\emph{bound} each other, can now be deduced from the two preceding, so that
the unity of reality and negation, identity and opposition, which makes up
\emph{the ground's} essential mark, already receives the form of proof. As
antithesis and synthesis thus condition and integrate each other in the
continued development, the philosophical categories take on an intermediate
character of fixity and fluidity, whereby the whole system rounds itself off
in an immanent demonstration. While Fichte thus heeded only thought's
necessity, and philosophized for philosophy's sake, Jacobi on the contrary
followed the promptings of the heart, and philosophized for the result's
sake. In the firm conviction that every philosophical demonstration must
consistently lead either to fatalism (Spinoza) or to atheism (Fichte), he
polemicized against the finite reflection's ``Durch= und Durch=Begriffe,''
and defended with rhetorical power reason's immediate cognition of God
against the abstract knowledge's nihilism. While Kant still acknowledged the
categories' subjective universal validity, only that they could not be
applied to the thing in itself, Jacobi, as Hegel has already remarked, went
a step further, and demonstrated the understanding-concepts' own finitude
and mechanical conditionedness. According to Kant, speculative science must
be given up, because it is once for all a psychological fact that the forms
of cognition lie in ourselves, and not in the essence of things; according
to Jacobi, true comprehension becomes impossible, because the concept, as
such, is itself finite, and consequently incommensurable with truth's
infinite content. If Kant attacked metaphysics chiefly from the side of the
stuff, without consummating the mechanical demonstration's own dissolution,
then Jacobi on the contrary attacked metaphysics' form, and sought thereby
to bring every philosophical \emph{method} into discredit.

\anm{4} When Kant calls the philosopher an \emph{artist of reason}, he
thereby intimates that philosophy cannot dispense with a determinate
structure, that its scientific content has its validity in the form; when
Jacobi on the contrary sets the philosopher the task of \emph{unveiling
existence} (Daseyn zu enthüllen), he thereby reminds us that science's true
worth always rests upon its substance. The philosophical question of method
is therefore as far from being an empty formal question as a mechanical
schema is from being a philosophical system. The form that determines the
content is in turn determined by the content's peculiar character; as
impossible as it is to develop a content without giving it form, just as
unscientific is it on the other side to confuse the content with the stuff,
and to apply an already finished form to a content given from without: that
which determines (the form) and that which is determined (the content)
reflect each other in the complete determinateness. He who would judge a
philosophy's substance without regard to its peculiar form proceeds from the
presupposition that the identity of content and form must rest upon the
difference of both. If this way of seeing had no validity, then the history
of philosophy would be meaningless; for the historical advance is always
conditioned by the inadequate and outdated forms being cast off, in order
that their essential content may be reproduced in a new form. He who
maintains that the philosophical form cannot be separated from philosophy's
content proceeds from the presupposition that the difference of form and
content grounds their true identity --- that the difference is thus a
moment vanishing in the identity. If this contention were altogether false,
science itself would be superfluous. --- The difference between life's and
philosophy's substance is in fact to be regarded as properly a difference of
form, since all that has truth in life must also be recognized as truth in
philosophy. What is peculiar in philosophy's content will thus rest upon its
peculiar form; the productive philosophers who have discovered new forms
have at the same time drawn forth new moments of content. --- What the
intellectual intuition is with respect to the content, the dialectical
reflection is with respect to the form. All speculation is intuition; the
philosophical vision apprehends things' fundamental relations, life's
shapes and existence's stages of development in large outlines; intuition is
thus a plastic vision, for which form and method give themselves immediately
together with content and substance. Without speculative genius there is no
philosophical productivity: philosophy cannot be learned (Schelling). So
long, however, as intuition has one-sided preponderance over reflection, the
scientific form is present only in its possibility; but even if philosophy's
genius has raised itself to ``the realm of the dawn,'' still the ideas'
genial breaking-forth and intuition's prophecies have significance only
insofar as philosophy is able to go out beyond the beginning, fulfil the
prophecy, and take up actuality's concrete content in articulated forms.
Philosophy is science, and distinguishes itself from poetry in this, that
its forms are won through independent reflection and an all-sidedly testing
meditation; it thus presupposes study, it justifies its method, and must so
far be capable of being learned (Hegel).

\anm{5} A philosophy's one-sidedness must not be confused with its
incompleteness. Incompleteness shows itself in important points being passed
over or only provisionally treated; but one-sidedness rests essentially upon
certain fundamental moments being pushed back, while others on the contrary
come forward with undue preponderance. A philosophy can be one-sided without
being incomplete, and conversely incomplete without being one-sided.
One-sidedness lies then in the systematic plan, in the method, and thus
refers at once both to content and to form. Spinoza's system is one-sided,
because substantiality here absorbs subjectivity, and even had Spinoza
developed the human concept of freedom ever so fully, he would still never
have been able to acknowledge freedom's reality without giving up his
principle. Fichte's idealism is one-sided, because it lets subjectivity
devour substantiality. True enough, Fichte has a place for the category of
substance; but even had he written folios on the concept of substantiality,
he would still, according to his fundamental view, never have recognized its
objective significance. The true philosophical method can therefore in
general be characterized as a totality-dialectic, which, to use Professor
Sibbern's expression, forms itself ``under an all-embracing conference and
debate of everything against everything'' (cfr.\ Sibbern: Om Philosophiens
Begreb, Natur og Væsen.\ Pag.\ 24; 81).

\divrule

% ---- Second Part, printed pp. 97--283 ----
% Printed p. 97 sets "anden Deel:" LOWERCASE (against p. 6's "Første Deel:"
% and the Indhold's "Anden Deel."); the transcription records both readings.
% The lowercase is mirrored here.
\deel{second Part:}{The Doctrine of the Objective Concept.}
\division{Introduction.}

\parag{12}{The Dialectical-Speculative Method.}

% [text to be added: pp. 97--113]

\divrule

% The Capitel arguments of the Anden Deel (pp. 114, 178, 247) ARE letterspaced
% in the original, unlike pp. 6/35/68 — hence \capitelsp, as in the transcription.
\capitelsp{First Chapter.}{The Formation of the Objective Concept.}
\parag{13}{Being and Essence.}

% [text to be added: pp. 114--131]

\parag{14}{Essence and Actuality.}

% [text to be added: pp. 132--158]

\parag{15}{Actuality and the Concept.}

% [text to be added: pp. 159--177]

\divrule

\capitelsp{Second Chapter.}{The Freedom of the Objective Concept.}
\parag{16}{Ideality and Reality.}

% [text to be added: pp. 178--202]

\parag{17}{Subjectivity and Objectivity.}

% [text to be added: pp. 203--224]

\parag{18}{Teleology.}

% [text to be added: pp. 225--246]

\divrule

\capitelsp{Third Chapter.}{The Fullness of the Objective Concept: the Logic of the Idea.}
\parag{19}{The Originality of the Idea.}

% [text to be added: pp. 247--257]

\parag{20}{The Interaction of the Ideas.}

% [text to be added: pp. 258--276]

\parag{21}{Idea and God.}

% [text to be added: pp. 277--283]

\divrule

\begin{center}\rule{0.35\linewidth}{0.4pt}\end{center}

% ============================================================
% ERRATA ("Rettelser") — the original's errata leaf, for the record.
% All nine corrections are APPLIED in the Danish transcription (and hence
% reflected in this translation); the leaf is reproduced in translation here.
% The p. 92 entry misprints "ci" for "ei" in its own right-hand column.
% ============================================================
\cleardoublepage
\chapter*{Errata}
\markboth{Errata}{Errata}
{\small\setlength{\parindent}{0pt}
Page 80 line 3 from above Omsætningen read: Oversætningen\par
--- \quad 92 \quad --- \quad 13 \quad --- \quad A er ei = A, \quad --- \quad A er ci = --- A\par
--- \quad 93 \quad --- \quad 7 \quad --- \quad Betingelser \quad --- \quad Betingethed\par
--- \quad 112 \quad --- \quad 4 from below \quad Daddel \quad --- \quad Dadel\par
--- \quad 152 \quad --- \quad 8 \quad --- \quad med \quad --- \quad ved\par
--- \quad 154 \quad --- \quad 15 \quad --- \quad Dyr=Rige \quad --- \quad Dyre=Rige\par
--- \quad 187 \quad --- \quad 9 \quad --- \quad εἵδη \quad --- \quad εἴδη\par
--- \quad 206 \quad --- \quad 9 \quad --- \quad en miniatur \quad --- \quad en miniature\par
--- \quad 241 \quad --- \quad 1 \quad --- \quad θετον \quad --- \quad θεῖον\par
}

\end{document}
